\section{动态规划进阶和容器源码章正文案例推导补强}

这一节补齐本章正文出现过的 LeetCode 案例推导。阅读顺序统一按“题目说明、样例入口、暴力基线、暴力过程、重复工作、优化条件、相邻状态、状态复用、指针和字段、代码落点、正确性、边界实验、复杂度变化、迁移追问”展开；目的不是新增题量，而是把正文中的案例都讲成可复述、可证明、可迁移的解题链。

\subsection{LeetCode 416 分割等和子集}

\textbf{题目说明。} LeetCode 416 分割等和子集 的题面入口要先还原成具体任务：总和一半是背包容量，每个数只能使用一次。

\textbf{样例入口。} 拿 [1, 5, 11, 5] 手算，总和 22，目标容量 11；可以用 11 单独凑成，也可以 5+5+1。

\textbf{暴力基线。} 暴力枚举每个数放到左集合还是右集合，最后检查两边和是否相等。

\textbf{暴力过程。} 以 [1, 5, 11, 5] 为例，总和 22，目标是找一个子集和为 11；可以选单独的 11，也可以选 5+5+1。

\textbf{重复工作。} 题面模型：总和一半是背包容量，每个数只能使用一次。暴力版的慢点要落到本题：浪费在不同选择顺序会反复得到相同剩余容量。问题只关心能否凑到 sum/2，不关心具体左右集合顺序。后面的优化只消掉这类重复工作，不能改变答案集合。

\textbf{优化条件。} 本题的关键观察是：0-1 背包可达性，一维容量必须倒序。这句话必须能翻译成可维护状态；如果题目条件改变后观察不再成立，就不能继续套原来的写法。

\textbf{相邻状态。} 规律是分割等和转成 0-1 背包是否能凑到 sum/2，容量必须倒序避免重复使用同一元素。把这个特例继续拆开看：阶段是物品，容量是资源，状态值是可达、最值还是计数。一个物品能否在同一轮再次使用，直接决定容量循环方向；组合和排列也由循环顺序表达。

\textbf{状态复用。} 先判断总和是否为偶数。dp[cap] 表示处理过若干数字后是否能凑出 cap，每个 num 只能用一次，所以容量倒序更新。手算时只改被新输入影响的那一小部分，而不是回头重算完整候选。

\textbf{指针和字段。} 状态字段要能说清：target=sum/2，cap 是容量，num 是当前物品；dp[cap-num] 为真则 dp[cap] 可变真。手算时按这个协议跟踪：处理 1 后可达 1；处理 5 后可达 5 和 6；处理 11 后 target 11 直接可达。

\textbf{代码落点。} 关键是 for cap from target down to num。若正序更新，会在同一轮把当前 num 用多次，变成完全背包。关键代码行必须能逐句翻译成“旧状态怎样变成新状态”，否则就只是模板。

\textbf{正确性。} 每个数只有选或不选两种状态；倒序容量保证转移读取的是上一阶段结果，不会重复使用当前数。

\textbf{边界实验。} 先用总和为奇数、目标为零、数值较大、复杂度依赖 target做反例检查。实验时覆盖总和为奇数、目标为零、数值较大、复杂度依赖 target，并故意把容量方向写反构造最小反例。这样能确认循环方向来自题意，不是背模板。

\textbf{复杂度变化。} 暴力 O(\ensuremath{2^{n}})，0-1 背包 O(n*target) 时间、O(target) 空间。

\textbf{迁移追问。} 如果题目改成“负数会破坏普通容量模型，需要重新建模”，先判断原状态是否仍然覆盖新答案集合，再决定是微调边界、改 value 语义，还是换算法。

\subsection{LeetCode 494 目标和}

\textbf{题目说明。} LeetCode 494 目标和 的题面入口要先还原成具体任务：正负号选择可代数转换为选若干数凑出特定正和。

\textbf{样例入口。} 拿 [1, 1, 1, 1, 1]、target=3 手算，给其中一个 1 加负号，其余为正可以得到 3，共有 5 种位置选择。也可转成选正数子集和为 (sum+target)/2。

\textbf{暴力基线。} 暴力递归给每个数添加 + 或 -，枚举 \ensuremath{2^{n}} 个符号序列并统计和为 target 的方案。

\textbf{暴力过程。} 以 [1, 1, 1, 1, 1]、target=3 为例，只有一个 1 取负、其余取正时可得到 3，共 5 种位置选择。

\textbf{重复工作。} 题面模型：正负号选择可代数转换为选若干数凑出特定正和。暴力版的慢点要落到本题：浪费在符号序列会反复得到相同部分和。把正数集合 P、负数集合 N 写成 sum(P)-sum(N)=target，可转成子集和计数。后面的优化只消掉这类重复工作，不能改变答案集合。

\textbf{优化条件。} 本题的关键观察是：若 S 加 target 为奇数或目标绝对值过大，直接无解。这句话必须能翻译成可维护状态；如果题目条件改变后观察不再成立，就不能继续套原来的写法。

\textbf{相邻状态。} 规律是符号分配变成 0-1 计数背包，奇偶和范围先判定。把这个特例继续拆开看：阶段是物品，容量是资源，状态值是可达、最值还是计数。一个物品能否在同一轮再次使用，直接决定容量循环方向；组合和排列也由循环顺序表达。

\textbf{状态复用。} 总和 total。若 total+target 为奇数或 abs(target)>total，直接 0；否则找子集和 positive=(total+target)/2 的方案数。手算时只改被新输入影响的那一小部分，而不是回头重算完整候选。

\textbf{指针和字段。} 状态字段要能说清：dp[cap] 是凑出 cap 的方案数，num 是当前数；0-1 计数背包容量倒序更新。手算时按这个协议跟踪：target=3、total=5，positive=4；选择四个 1 作为正数等价于选择哪个 1 不放进去，共 5 种。

\textbf{代码落点。} 关键是 dp[0]=1，for cap from positive down to num。数字 0 会让方案数翻倍，倒序更新能自然处理。关键代码行必须能逐句翻译成“旧状态怎样变成新状态”，否则就只是模板。

\textbf{正确性。} 每个符号分配和一个正数子集一一对应；子集和满足 positive 时，正负差正好为 target。

\textbf{边界实验。} 先用零元素会让方案数翻倍、目标为负、容量为零做反例检查。实验时覆盖零元素会让方案数翻倍、目标为负、容量为零，并故意把容量方向写反构造最小反例。这样能确认循环方向来自题意，不是背模板。

\textbf{复杂度变化。} 暴力 O(\ensuremath{2^{n}})，背包 O(n*positive) 时间，O(positive) 空间。

\textbf{迁移追问。} 如果题目改成“也可用哈希 DP 保存所有可能和，但背包转换更高效”，先判断原状态是否仍然覆盖新答案集合，再决定是微调边界、改 value 语义，还是换算法。

\subsection{LeetCode 518 零钱兑换 II}

\textbf{题目说明。} LeetCode 518 零钱兑换 II 的题面入口要先还原成具体任务：硬币无限使用，问组合数而不是最少数量。

\textbf{样例入口。} 拿 amount=5、coins=[1, 2, 5] 手算，组合 5、2+2+1、2+1+1+1、全 1 一共 4 种；1+2 和 2+1 不能重复算。

\textbf{暴力基线。} 暴力 DFS 枚举每种硬币使用几个，或者枚举所有硬币排列再去重。

\textbf{暴力过程。} amount=5、coins=[1, 2, 5] 时，组合有 5、2+2+1、2+1+1+1、1+1+1+1+1，共 4 种；1+2 和 2+1 不能重复算。

\textbf{重复工作。} 题面模型：硬币无限使用，问组合数而不是最少数量。暴力版的慢点要落到本题：浪费在排列顺序重复。组合计数应固定硬币顺序，让每种组合只在一种顺序下被统计。后面的优化只消掉这类重复工作，不能改变答案集合。

\textbf{优化条件。} 本题的关键观察是：外层遍历硬币、内层金额正序，保证组合按固定硬币顺序统计。这句话必须能翻译成可维护状态；如果题目条件改变后观察不再成立，就不能继续套原来的写法。

\textbf{相邻状态。} 规律是组合计数时外层枚举硬币，内层容量正序。把这个特例继续拆开看：阶段是物品，容量是资源，状态值是可达、最值还是计数。一个物品能否在同一轮再次使用，直接决定容量循环方向；组合和排列也由循环顺序表达。

\textbf{状态复用。} 完全背包计数。dp[amount] 表示当前已处理硬币集合能凑出该金额的组合数；外层枚举 coin，内层金额正序。手算时只改被新输入影响的那一小部分，而不是回头重算完整候选。

\textbf{指针和字段。} 状态字段要能说清：coin 是当前允许使用的硬币，cap 是当前金额，dp[cap-coin] 是加入一枚 coin 前的组合数。手算时按这个协议跟踪：处理 coin=2 时，dp[4] 会从 dp[2] 来，表示使用两个 2；不会再把 1+2 和 2+1 分开统计。

\textbf{代码落点。} 关键循环是 for coin in coins: for cap from coin to amount: dp[cap]+=dp[cap-coin]。若外层金额内层硬币，会统计排列数。关键代码行必须能逐句翻译成“旧状态怎样变成新状态”，否则就只是模板。

\textbf{正确性。} 外层固定硬币顺序后，每个组合按最大硬币处理阶段唯一计入；金额正序允许当前硬币重复使用。

\textbf{边界实验。} 先用amount 为零、无硬币、组合数较大、循环顺序反例做反例检查。实验时覆盖amount 为零、无硬币、组合数较大、循环顺序反例，并故意把容量方向写反构造最小反例。这样能确认循环方向来自题意，不是背模板。

\textbf{复杂度变化。} 暴力指数级；完全背包计数 O(amount*coins.size()) 时间，O(amount) 空间。

\textbf{迁移追问。} 如果题目改成“若外层金额内层硬币，会统计排列数”，先判断原状态是否仍然覆盖新答案集合，再决定是微调边界、改 value 语义，还是换算法。

\subsection{LeetCode 121 买卖股票的最佳时机}

\textbf{题目说明。} LeetCode 121 买卖股票的最佳时机 的题面入口要先还原成具体任务：卖出日固定时，最佳买入日是此前最低价格。

\textbf{样例入口。} 拿价格 [7, 1, 5] 手算。站在 7 不能卖出获利；站在 1 更新最低买入价；站在 5 时利润是 5-1。

\textbf{暴力基线。} 暴力枚举买入日 i 和卖出日 j，要求 i<j，计算 prices[j]-prices[i] 的最大值。

\textbf{暴力过程。} 以 [7, 1, 5, 3, 6, 4] 为例，若卖出日在价格 6，最好的买入日是此前最低价 1，利润 5。

\textbf{重复工作。} 题面模型：卖出日固定时，最佳买入日是此前最低价格。暴力版的慢点要落到本题：浪费在每个卖出日都回头找最低买入价。扫描过程中只需要维护历史最低价。后面的优化只消掉这类重复工作，不能改变答案集合。

\textbf{优化条件。} 本题的关键观察是：扫描维护历史最低价和最大利润，一次交易无需复杂状态机。这句话必须能翻译成可维护状态；如果题目条件改变后观察不再成立，就不能继续套原来的写法。

\textbf{相邻状态。} 规律是固定卖出日后，最佳买入日只需要历史最低价，状态不是所有买卖对。把这个特例继续拆开看：先问暴力搜索里哪些不同路径到达同一个后续问题，再给这个后续问题命名为状态。状态必须包含未来决策需要的全部信息，转移只从更小且已经算清的状态来。

\textbf{状态复用。} min\_price 保存当前日前的最低价格，best 保存最大利润。每天先用当前价格尝试卖出，再更新 min\_price。手算时只改被新输入影响的那一小部分，而不是回头重算完整候选。

\textbf{指针和字段。} 状态字段要能说清：price 是当前卖出价候选，min\_price 是历史最优买入价，best 是已见最大利润。手算时按这个协议跟踪：扫描到 5 时 min\_price=1，best=4；扫描到 6 时利润 5，更新 best。

\textbf{代码落点。} 核心顺序是先用当前价格尝试卖出更新 best，再用当前价格更新 min\_price。若先更新最低价，同一天买卖利润为 0，不影响答案但语义要讲清楚。关键代码行必须能逐句翻译成“旧状态怎样变成新状态”，否则就只是模板。

\textbf{正确性。} 任意最优交易按卖出日分类；固定卖出日时，最佳买入日一定是它之前最低价格，状态完整覆盖这个选择。

\textbf{边界实验。} 先用单日价格、一直下跌、一直上涨、价格相同做反例检查。实验时覆盖单日价格、一直下跌、一直上涨、价格相同，先打印完整 DP 表，再比较空间压缩版本。若压缩版不同，优先检查遍历顺序和旧值保存。

\textbf{复杂度变化。} 暴力 O(\ensuremath{n^{2}})，一次扫描 O(n) 时间、O(1) 空间。

\textbf{迁移追问。} 如果题目改成“多次交易、冷冻期和手续费版本都要扩展成状态机”，先判断原状态是否仍然覆盖新答案集合，再决定是微调边界、改 value 语义，还是换算法。

\subsection{LeetCode 309 最佳买卖股票含冷冻期}

\textbf{题目说明。} LeetCode 309 最佳买卖股票含冷冻期 的题面入口要先还原成具体任务：冷冻期让买入来源受昨天卖出状态限制。

\textbf{样例入口。} 拿 [1, 2, 3, 0, 2] 手算，第 2 天卖出后第 3 天不能立刻买入，因为有冷冻期；状态必须区分持股、刚卖出、空仓可买。

\textbf{暴力基线。} 慢版本先写递归搜索树：节点表示处理位置、剩余资源和当前选择。只要不同路径反复到达同一个节点，就说明需要把这个节点命名成 DP 状态。放到本题就是：先按“冷冻期让买入来源受昨天卖出状态限制”完整试慢版本；样例里已经能看到“规律是冷冻期把普通买卖股票拆成多个日末状态，转移不能跨过限制”，所以优化前必须先能说清慢版本枚举了哪些候选。

\textbf{暴力过程。} 拿 [1, 2, 3, 0, 2] 手算，第 2 天卖出后第 3 天不能立刻买入，因为有冷冻期；状态必须区分持股、刚卖出、空仓可买。

\textbf{重复工作。} 题面模型：冷冻期让买入来源受昨天卖出状态限制。暴力版的慢点要落到本题：慢版本会在“冷冻期让买入来源受昨天卖出状态限制”的选择树里反复到达相同参数。样例规律是“规律是冷冻期把普通买卖股票拆成多个日末状态，转移不能跨过限制”，说明这个参数组合可以命名成状态，算过之后后面直接复用。后面的优化只消掉这类重复工作，不能改变答案集合。

\textbf{优化条件。} 本题的关键观察是：状态机维护持有、刚卖出、休息三种日终状态。这句话必须能翻译成可维护状态；如果题目条件改变后观察不再成立，就不能继续套原来的写法。

\textbf{相邻状态。} 规律是冷冻期把普通买卖股票拆成多个日末状态，转移不能跨过限制。把这个特例继续拆开看：先问暴力搜索里哪些不同路径到达同一个后续问题，再给这个后续问题命名为状态。状态必须包含未来决策需要的全部信息，转移只从更小且已经算清的状态来。

\textbf{状态复用。} 把重复递归节点命名为状态；遍历顺序只是在保证依赖状态已经算完。本题落点是：规律是冷冻期把普通买卖股票拆成多个日末状态，转移不能跨过限制。手算时只改被新输入影响的那一小部分，而不是回头重算完整候选。

\textbf{指针和字段。} 状态字段要能说清：状态下标、状态值语义、初始化、转移来源、遍历顺序；空间压缩时还要指出读到的是旧值还是新值。手算时按这个协议跟踪：拿 [1, 2, 3, 0, 2] 手算，第 2 天卖出后第 3 天不能立刻买入，因为有冷冻期；状态必须区分持股、刚卖出、空仓可买。然后按协议复查：手算时先写一个很小的 DP 表或记忆化搜索记录。每个格子旁边写它来自哪些更小状态。

\textbf{代码落点。} 先写未压缩版本校验转移；压缩前后用小表对拍；不可达哨兵参与加法前要检查溢出。关键代码行必须能逐句翻译成“旧状态怎样变成新状态”，否则就只是模板。

\textbf{正确性。} 证明从最后一步开始：任意最优解的最后一步都在转移枚举中，转移来源已经由更小状态正确计算。

\textbf{边界实验。} 先用空价格、一天、连续上涨、卖出后不能次日买入做反例检查。实验时覆盖空价格、一天、连续上涨、卖出后不能次日买入，先打印完整 DP 表，再比较空间压缩版本。若压缩版不同，优先检查遍历顺序和旧值保存。

\textbf{复杂度变化。} 暴力递归会重复计算相同子问题，常见指数级；DP 把每个状态算一次，时间是状态数乘单个转移成本，空间是状态表大小或压缩后的大小。

\textbf{迁移追问。} 如果题目改成“手续费版本可以在买入或卖出转移中扣费”，先判断原状态是否仍然覆盖新答案集合，再决定是微调边界、改 value 语义，还是换算法。

\subsection{LeetCode 714 买卖股票含手续费}

\textbf{题目说明。} LeetCode 714 买卖股票含手续费 的题面入口要先还原成具体任务：手续费改变交易收益，仍可用持有和不持有状态。

\textbf{样例入口。} 拿 [1, 3, 2, 8, 4, 9]、fee=2 手算，8 卖出时净赚 5；手续费让频繁小波动不一定值得交易。

\textbf{暴力基线。} 慢版本先写递归搜索树：节点表示处理位置、剩余资源和当前选择。只要不同路径反复到达同一个节点，就说明需要把这个节点命名成 DP 状态。放到本题就是：先按“手续费改变交易收益，仍可用持有和不持有状态”完整试慢版本；样例里已经能看到“规律是每天维护 hold 和 cash，卖出时扣手续费，买入时从 cash 转到 hold”，所以优化前必须先能说清慢版本枚举了哪些候选。

\textbf{暴力过程。} 拿 [1, 3, 2, 8, 4, 9]、fee=2 手算，8 卖出时净赚 5；手续费让频繁小波动不一定值得交易。

\textbf{重复工作。} 题面模型：手续费改变交易收益，仍可用持有和不持有状态。暴力版的慢点要落到本题：慢版本会在“手续费改变交易收益，仍可用持有和不持有状态”的选择树里反复到达相同参数。样例规律是“规律是每天维护 hold 和 cash，卖出时扣手续费，买入时从 cash 转到 hold”，说明这个参数组合可以命名成状态，算过之后后面直接复用。后面的优化只消掉这类重复工作，不能改变答案集合。

\textbf{优化条件。} 本题的关键观察是：卖出时扣费或买入时扣费都可以，但不要重复扣。这句话必须能翻译成可维护状态；如果题目条件改变后观察不再成立，就不能继续套原来的写法。

\textbf{相邻状态。} 规律是每天维护 hold 和 cash，卖出时扣手续费，买入时从 cash 转到 hold。把这个特例继续拆开看：先问暴力搜索里哪些不同路径到达同一个后续问题，再给这个后续问题命名为状态。状态必须包含未来决策需要的全部信息，转移只从更小且已经算清的状态来。

\textbf{状态复用。} 把重复递归节点命名为状态；遍历顺序只是在保证依赖状态已经算完。本题落点是：规律是每天维护 hold 和 cash，卖出时扣手续费，买入时从 cash 转到 hold。手算时只改被新输入影响的那一小部分，而不是回头重算完整候选。

\textbf{指针和字段。} 状态字段要能说清：状态下标、状态值语义、初始化、转移来源、遍历顺序；空间压缩时还要指出读到的是旧值还是新值。手算时按这个协议跟踪：拿 [1, 3, 2, 8, 4, 9]、fee=2 手算，8 卖出时净赚 5；手续费让频繁小波动不一定值得交易。然后按协议复查：手算时先写一个很小的 DP 表或记忆化搜索记录。每个格子旁边写它来自哪些更小状态。

\textbf{代码落点。} 先写未压缩版本校验转移；压缩前后用小表对拍；不可达哨兵参与加法前要检查溢出。关键代码行必须能逐句翻译成“旧状态怎样变成新状态”，否则就只是模板。

\textbf{正确性。} 证明从最后一步开始：任意最优解的最后一步都在转移枚举中，转移来源已经由更小状态正确计算。

\textbf{边界实验。} 先用手续费为零、价格震荡、小利润不足手续费做反例检查。实验时覆盖手续费为零、价格震荡、小利润不足手续费，先打印完整 DP 表，再比较空间压缩版本。若压缩版不同，优先检查遍历顺序和旧值保存。

\textbf{复杂度变化。} 暴力递归会重复计算相同子问题，常见指数级；DP 把每个状态算一次，时间是状态数乘单个转移成本，空间是状态表大小或压缩后的大小。

\textbf{迁移追问。} 如果题目改成“这是状态机比局部贪心更稳的例子”，先判断原状态是否仍然覆盖新答案集合，再决定是微调边界、改 value 语义，还是换算法。

\subsection{LeetCode 188 买卖股票 IV}

\textbf{题目说明。} LeetCode 188 买卖股票 IV 的题面入口要先还原成具体任务：最多 k 次交易让状态必须同时包含交易次数和是否持有股票。

\textbf{样例入口。} 拿 k=2、价格 [3, 2, 6] 手算。某天结束时只说最大利润不够，因为未来能否卖出取决于是否持股，也取决于已经完成几次交易。

\textbf{暴力基线。} 慢版本先写递归搜索树：节点表示处理位置、剩余资源和当前选择。只要不同路径反复到达同一个节点，就说明需要把这个节点命名成 DP 状态。放到本题就是：先按“最多 k 次交易让状态必须同时包含交易次数和是否持有股票”完整试慢版本；样例里已经能看到“规律是状态必须包含交易次数和持有状态；k 很大时才退化为无限次交易”，所以优化前必须先能说清慢版本枚举了哪些候选。

\textbf{暴力过程。} 拿 k=2、价格 [3, 2, 6] 手算。某天结束时只说最大利润不够，因为未来能否卖出取决于是否持股，也取决于已经完成几次交易。

\textbf{重复工作。} 题面模型：最多 k 次交易让状态必须同时包含交易次数和是否持有股票。暴力版的慢点要落到本题：慢版本会在“最多 k 次交易让状态必须同时包含交易次数和是否持有股票”的选择树里反复到达相同参数。样例规律是“规律是状态必须包含交易次数和持有状态；k 很大时才退化为无限次交易”，说明这个参数组合可以命名成状态，算过之后后面直接复用。后面的优化只消掉这类重复工作，不能改变答案集合。

\textbf{优化条件。} 本题的关键观察是：维护每个交易次数下的 buy 和 sell 状态；当 k 足够大时可退化为无限次交易。这句话必须能翻译成可维护状态；如果题目条件改变后观察不再成立，就不能继续套原来的写法。

\textbf{相邻状态。} 规律是状态必须包含交易次数和持有状态；k 很大时才退化为无限次交易。把这个特例继续拆开看：先问暴力搜索里哪些不同路径到达同一个后续问题，再给这个后续问题命名为状态。状态必须包含未来决策需要的全部信息，转移只从更小且已经算清的状态来。

\textbf{状态复用。} 把重复递归节点命名为状态；遍历顺序只是在保证依赖状态已经算完。本题落点是：规律是状态必须包含交易次数和持有状态；k 很大时才退化为无限次交易。手算时只改被新输入影响的那一小部分，而不是回头重算完整候选。

\textbf{指针和字段。} 状态字段要能说清：状态下标、状态值语义、初始化、转移来源、遍历顺序；空间压缩时还要指出读到的是旧值还是新值。手算时按这个协议跟踪：拿 k=2、价格 [3, 2, 6] 手算。某天结束时只说最大利润不够，因为未来能否卖出取决于是否持股，也取决于已经完成几次交易。然后按协议复查：手算时先写一个很小的 DP 表或记忆化搜索记录。每个格子旁边写它来自哪些更小状态。

\textbf{代码落点。} 先写未压缩版本校验转移；压缩前后用小表对拍；不可达哨兵参与加法前要检查溢出。关键代码行必须能逐句翻译成“旧状态怎样变成新状态”，否则就只是模板。

\textbf{正确性。} 证明从最后一步开始：任意最优解的最后一步都在转移枚举中，转移来源已经由更小状态正确计算。

\textbf{边界实验。} 先用k 为零、价格天数很少、k 大于等于 n/2、更新顺序造成同一天状态污染做反例检查。实验时覆盖k 为零、价格天数很少、k 大于等于 n/2、更新顺序造成同一天状态污染，先打印完整 DP 表，再比较空间压缩版本。若压缩版不同，优先检查遍历顺序和旧值保存。

\textbf{复杂度变化。} 暴力递归会重复计算相同子问题，常见指数级；DP 把每个状态算一次，时间是状态数乘单个转移成本，空间是状态表大小或压缩后的大小。

\textbf{迁移追问。} 如果题目改成“手续费、冷冻期和最多两次交易都是同一状态机框架的不同约束”，先判断原状态是否仍然覆盖新答案集合，再决定是微调边界、改 value 语义，还是换算法。

\subsection{LeetCode 300 最长递增子序列}

\textbf{题目说明。} LeetCode 300 最长递增子序列 的题面入口要先还原成具体任务：二次 DP 固定结尾，二分优化维护每个长度的最小结尾。

\textbf{样例入口。} 拿 [10, 9, 2, 5, 3] 手算。二次 DP 固定每个位置为结尾，5 可接在 2 后；优化时 tails[长度] 保存该长度递增子序列的最小尾值，尾值越小未来越容易接。

\textbf{暴力基线。} 暴力递归或 DP 枚举每个元素作为子序列结尾，向前寻找所有更小元素。

\textbf{暴力过程。} 以 [10, 9, 2, 5, 3, 7, 101, 18] 为例，二次 DP 可以算出以 7 结尾的长度来自 2、5 或 3 后面接 7。

\textbf{重复工作。} 题面模型：二次 DP 固定结尾，二分优化维护每个长度的最小结尾。暴力版的慢点要落到本题：二次 DP 慢在每个位置都回看所有历史元素。更快写法不需要保存真实序列，只保存每个长度的最小可能结尾。后面的优化只消掉这类重复工作，不能改变答案集合。

\textbf{优化条件。} 本题的关键观察是：tails 不是实际答案序列，而是未来扩展潜力表。这句话必须能翻译成可维护状态；如果题目条件改变后观察不再成立，就不能继续套原来的写法。

\textbf{相邻状态。} 规律是二分替换 tails，不代表真实序列完整内容，而是维护未来潜力。把这个特例继续拆开看：先问暴力搜索里哪些不同路径到达同一个后续问题，再给这个后续问题命名为状态。状态必须包含未来决策需要的全部信息，转移只从更小且已经算清的状态来。

\textbf{状态复用。} tails[len] 表示长度为 len+1 的递增子序列中最小结尾。对每个 x，用 lower\_bound 找第一个 >=x 的位置并替换。手算时只改被新输入影响的那一小部分，而不是回头重算完整候选。

\textbf{指针和字段。} 状态字段要能说清：tails 不是最终答案序列，而是未来扩展潜力表；tails.size() 是当前最长长度。手算时按这个协议跟踪：处理 2、5、3 时，tails 从 [2, 5] 变成 [2, 3]，长度没变，但结尾更小，未来更容易接 7。

\textbf{代码落点。} 严格递增用 lower\_bound；若题目改成非递减，要改用 upper\_bound。要还原路径需要额外 predecessor 数组。关键代码行必须能逐句翻译成“旧状态怎样变成新状态”，否则就只是模板。

\textbf{正确性。} 同样长度下，结尾越小，未来能接上的数字集合只会更多；替换 tails 不改变已知长度，只提升扩展潜力。

\textbf{边界实验。} 先用重复元素、严格递增和非递减区别、空数组做反例检查。实验时覆盖重复元素、严格递增和非递减区别、空数组，先打印完整 DP 表，再比较空间压缩版本。若压缩版不同，优先检查遍历顺序和旧值保存。

\textbf{复杂度变化。} 二次 DP O(\ensuremath{n^{2}})，tails 加二分 O(n log n) 时间、O(n) 空间。

\textbf{迁移追问。} 如果题目改成“若要还原路径，需要额外前驱数组，不能只看 tails”，先判断原状态是否仍然覆盖新答案集合，再决定是微调边界、改 value 语义，还是换算法。

\subsection{LeetCode 62 不同路径}

\textbf{题目说明。} LeetCode 62 不同路径 的题面入口要先还原成具体任务：网格状态由上方和左方两种最后一步转移而来。

\textbf{样例入口。} 拿 2x3 网格手算，到每个格子的最后一步只能从上方或左方来；第一行和第一列都只有 1 条路，右下角等于左边 2 加上上边 1，得到 3。

\textbf{暴力基线。} 慢版本先写递归搜索树：节点表示处理位置、剩余资源和当前选择。只要不同路径反复到达同一个节点，就说明需要把这个节点命名成 DP 状态。放到本题就是：先按“网格状态由上方和左方两种最后一步转移而来”完整试慢版本；样例里已经能看到“规律是 dp[row][col]=上方+左方，滚动数组只是在保存上一行同列和当前行左侧”，所以优化前必须先能说清慢版本枚举了哪些候选。

\textbf{暴力过程。} 拿 2x3 网格手算，到每个格子的最后一步只能从上方或左方来；第一行和第一列都只有 1 条路，右下角等于左边 2 加上上边 1，得到 3。

\textbf{重复工作。} 题面模型：网格状态由上方和左方两种最后一步转移而来。暴力版的慢点要落到本题：慢版本会在“网格状态由上方和左方两种最后一步转移而来”的选择树里反复到达相同参数。样例规律是“规律是 dp[row][col]=上方+左方，滚动数组只是在保存上一行同列和当前行左侧”，说明这个参数组合可以命名成状态，算过之后后面直接复用。后面的优化只消掉这类重复工作，不能改变答案集合。

\textbf{优化条件。} 本题的关键观察是：二维表可压缩为一维，因为当前行只依赖上一行同列和当前行左侧。这句话必须能翻译成可维护状态；如果题目条件改变后观察不再成立，就不能继续套原来的写法。

\textbf{相邻状态。} 规律是 dp[row][col]=上方+左方，滚动数组只是在保存上一行同列和当前行左侧。把这个特例继续拆开看：先问暴力搜索里哪些不同路径到达同一个后续问题，再给这个后续问题命名为状态。状态必须包含未来决策需要的全部信息，转移只从更小且已经算清的状态来。

\textbf{状态复用。} 把重复递归节点命名为状态；遍历顺序只是在保证依赖状态已经算完。本题落点是：规律是 dp[row][col]=上方+左方，滚动数组只是在保存上一行同列和当前行左侧。手算时只改被新输入影响的那一小部分，而不是回头重算完整候选。

\textbf{指针和字段。} 状态字段要能说清：状态下标、状态值语义、初始化、转移来源、遍历顺序；空间压缩时还要指出读到的是旧值还是新值。手算时按这个协议跟踪：拿 2x3 网格手算，到每个格子的最后一步只能从上方或左方来；第一行和第一列都只有 1 条路，右下角等于左边 2 加上上边 1，得到 3。然后按协议复查：手算时先写一个很小的 DP 表或记忆化搜索记录。每个格子旁边写它来自哪些更小状态。

\textbf{代码落点。} 先写未压缩版本校验转移；压缩前后用小表对拍；不可达哨兵参与加法前要检查溢出。关键代码行必须能逐句翻译成“旧状态怎样变成新状态”，否则就只是模板。

\textbf{正确性。} 证明从最后一步开始：任意最优解的最后一步都在转移枚举中，转移来源已经由更小状态正确计算。

\textbf{边界实验。} 先用一行、一列、较大网格结果溢出、障碍物版本边界做反例检查。实验时覆盖一行、一列、较大网格结果溢出、障碍物版本边界，先打印完整 DP 表，再比较空间压缩版本。若压缩版不同，优先检查遍历顺序和旧值保存。

\textbf{复杂度变化。} 暴力递归会重复计算相同子问题，常见指数级；DP 把每个状态算一次，时间是状态数乘单个转移成本，空间是状态表大小或压缩后的大小。

\textbf{迁移追问。} 如果题目改成“带障碍物时遇到障碍状态清零，第一行第一列要按阻断处理”，先判断原状态是否仍然覆盖新答案集合，再决定是微调边界、改 value 语义，还是换算法。

\subsection{LeetCode 64 最小路径和}

\textbf{题目说明。} LeetCode 64 最小路径和 的题面入口要先还原成具体任务：状态是到达当前格子的最小代价，最后一步来自上方或左方。

\textbf{样例入口。} 拿 [[1, 3, 1], [1, 5, 1], [4, 2, 1]] 手算，到 5 的最小代价来自左侧 4 或上方 4，再加 5；右下角只需要比较上方和左方的已知最小代价。

\textbf{暴力基线。} 慢版本先写递归搜索树：节点表示处理位置、剩余资源和当前选择。只要不同路径反复到达同一个节点，就说明需要把这个节点命名成 DP 状态。放到本题就是：先按“状态是到达当前格子的最小代价，最后一步来自上方或左方”完整试慢版本；样例里已经能看到“规律是最后一步来自上或左，先初始化第一行第一列，内部统一取 min”，所以优化前必须先能说清慢版本枚举了哪些候选。

\textbf{暴力过程。} 拿 [[1, 3, 1], [1, 5, 1], [4, 2, 1]] 手算，到 5 的最小代价来自左侧 4 或上方 4，再加 5；右下角只需要比较上方和左方的已知最小代价。

\textbf{重复工作。} 题面模型：状态是到达当前格子的最小代价，最后一步来自上方或左方。暴力版的慢点要落到本题：慢版本会在“状态是到达当前格子的最小代价，最后一步来自上方或左方”的选择树里反复到达相同参数。样例规律是“规律是最后一步来自上或左，先初始化第一行第一列，内部统一取 min”，说明这个参数组合可以命名成状态，算过之后后面直接复用。后面的优化只消掉这类重复工作，不能改变答案集合。

\textbf{优化条件。} 本题的关键观察是：先初始化第一行和第一列，再处理内部，主转移保持简单。这句话必须能翻译成可维护状态；如果题目条件改变后观察不再成立，就不能继续套原来的写法。

\textbf{相邻状态。} 规律是最后一步来自上或左，先初始化第一行第一列，内部统一取 min。把这个特例继续拆开看：先问暴力搜索里哪些不同路径到达同一个后续问题，再给这个后续问题命名为状态。状态必须包含未来决策需要的全部信息，转移只从更小且已经算清的状态来。

\textbf{状态复用。} 把重复递归节点命名为状态；遍历顺序只是在保证依赖状态已经算完。本题落点是：规律是最后一步来自上或左，先初始化第一行第一列，内部统一取 min。手算时只改被新输入影响的那一小部分，而不是回头重算完整候选。

\textbf{指针和字段。} 状态字段要能说清：状态下标、状态值语义、初始化、转移来源、遍历顺序；空间压缩时还要指出读到的是旧值还是新值。手算时按这个协议跟踪：拿 [[1, 3, 1], [1, 5, 1], [4, 2, 1]] 手算，到 5 的最小代价来自左侧 4 或上方 4，再加 5；右下角只需要比较上方和左方的已知最小代价。然后按协议复查：手算时先写一个很小的 DP 表或记忆化搜索记录。每个格子旁边写它来自哪些更小状态。

\textbf{代码落点。} 先写未压缩版本校验转移；压缩前后用小表对拍；不可达哨兵参与加法前要检查溢出。关键代码行必须能逐句翻译成“旧状态怎样变成新状态”，否则就只是模板。

\textbf{正确性。} 证明从最后一步开始：任意最优解的最后一步都在转移枚举中，转移来源已经由更小状态正确计算。

\textbf{边界实验。} 先用单行、单列、权值为零、路径和溢出做反例检查。实验时覆盖单行、单列、权值为零、路径和溢出，先打印完整 DP 表，再比较空间压缩版本。若压缩版不同，优先检查遍历顺序和旧值保存。

\textbf{复杂度变化。} 暴力递归会重复计算相同子问题，常见指数级；DP 把每个状态算一次，时间是状态数乘单个转移成本，空间是状态表大小或压缩后的大小。

\textbf{迁移追问。} 如果题目改成“若允许四方向移动，变成最短路而不是普通网格 DP”，先判断原状态是否仍然覆盖新答案集合，再决定是微调边界、改 value 语义，还是换算法。

\subsection{LeetCode 72 编辑距离}

\textbf{题目说明。} LeetCode 72 编辑距离 的题面入口要先还原成具体任务：二维状态表示两个前缀之间的最少编辑操作。

\textbf{样例入口。} 拿 word1=ab、word2=ac 手算最后一个字符。若 b 改成 c，是 dp[1][1]+1；若删掉 b，是 dp[1][2]+1；若插入 c，是 dp[2][1]+1。字符相等时才直接继承左上。

\textbf{暴力基线。} 暴力递归比较两个字符串后缀：相等就同时跳过，不等就尝试插入、删除、替换三种操作。

\textbf{暴力过程。} 以 word1=ab、word2=ac 为例，末尾 b 和 c 不同。替换 b 成 c 来自 dp[1][1]+1；删除 b 来自 dp[1][2]+1；插入 c 来自 dp[2][1]+1。

\textbf{重复工作。} 题面模型：二维状态表示两个前缀之间的最少编辑操作。暴力版的慢点要落到本题：浪费在不同编辑序列会反复到达同一对前缀长度。只要 i 和 j 相同，后续最少编辑次数就相同。后面的优化只消掉这类重复工作，不能改变答案集合。

\textbf{优化条件。} 本题的关键观察是：末尾字符相同继承左上；不同则枚举插入、删除、替换。这句话必须能翻译成可维护状态；如果题目条件改变后观察不再成立，就不能继续套原来的写法。

\textbf{相邻状态。} 规律是每个格子枚举最后一次编辑操作，而不是凭感觉填三邻居最小。把这个特例继续拆开看：先问暴力搜索里哪些不同路径到达同一个后续问题，再给这个后续问题命名为状态。状态必须包含未来决策需要的全部信息，转移只从更小且已经算清的状态来。

\textbf{状态复用。} dp[i][j] 表示 word1 前 i 个字符变成 word2 前 j 个字符的最少操作数。按前缀长度从小到大填表。手算时只改被新输入影响的那一小部分，而不是回头重算完整候选。

\textbf{指针和字段。} 状态字段要能说清：i/j 是前缀长度；dp[i-1][j] 对应删除，dp[i][j-1] 对应插入，dp[i-1][j-1] 对应替换或字符相等继承。手算时按这个协议跟踪：空串到 ac 需要插入 2 次；ab 到空串需要删除 2 次。内部格子按末尾字符是否相等决定转移。

\textbf{代码落点。} 关键是第一行 dp[0][j]=j、第一列 dp[i][0]=i。字符相等时不要额外加一。关键代码行必须能逐句翻译成“旧状态怎样变成新状态”，否则就只是模板。

\textbf{正确性。} 任意最优编辑序列的最后一步只有插入、删除、替换或不操作四类，转移枚举了全部最后一步。

\textbf{边界实验。} 先用空串、完全相同、完全不同、操作代价不同做反例检查。实验时覆盖空串、完全相同、完全不同、操作代价不同，先打印完整 DP 表，再比较空间压缩版本。若压缩版不同，优先检查遍历顺序和旧值保存。

\textbf{复杂度变化。} 暴力递归指数级；二维 DP O(mn) 时间和空间，可滚动到 O(min(m, n)) 空间。

\textbf{迁移追问。} 如果题目改成“若只允许插入删除，问题退化到和最长公共子序列相关”，先判断原状态是否仍然覆盖新答案集合，再决定是微调边界、改 value 语义，还是换算法。

\subsection{LeetCode 1143 最长公共子序列}

\textbf{题目说明。} LeetCode 1143 最长公共子序列 的题面入口要先还原成具体任务：状态表示两个前缀的最长公共子序列长度。

\textbf{样例入口。} 拿 abcde 和 ace 手算，末尾 e 相等时答案来自两个前缀 abcd 和 ac 加一；若末尾不同，就取删掉一侧末尾后的较大值。

\textbf{暴力基线。} 暴力递归枚举两个字符串的子序列选择，或从末尾开始相等就选、不等就分支删一边。

\textbf{暴力过程。} 以 text1=abcde、text2=ace 为例，末尾 e 相等时，答案来自 abcd 和 ac 的 LCS 加一；若末尾不同，就只能丢掉一侧末尾试更优。

\textbf{重复工作。} 题面模型：状态表示两个前缀的最长公共子序列长度。暴力版的慢点要落到本题：浪费在很多前缀对会被重复计算。只要两个前缀长度 i、j 固定，它们的 LCS 长度就是固定状态。后面的优化只消掉这类重复工作，不能改变答案集合。

\textbf{优化条件。} 本题的关键观察是：末尾相同取左上加一，不同取上方和左方最大。这句话必须能翻译成可维护状态；如果题目条件改变后观察不再成立，就不能继续套原来的写法。

\textbf{相邻状态。} 规律是 dp[i][j] 表示两个前缀的 LCS 长度，转移只看左、上、左上。把这个特例继续拆开看：先问暴力搜索里哪些不同路径到达同一个后续问题，再给这个后续问题命名为状态。状态必须包含未来决策需要的全部信息，转移只从更小且已经算清的状态来。

\textbf{状态复用。} dp[i][j] 表示 text1 前 i 个字符和 text2 前 j 个字符的 LCS 长度。相等看左上加一，不等取上方和左方最大。手算时只改被新输入影响的那一小部分，而不是回头重算完整候选。

\textbf{指针和字段。} 状态字段要能说清：i/j 是前缀长度，dp[i-1][j] 表示丢 text1 末尾，dp[i][j-1] 表示丢 text2 末尾，dp[i-1][j-1] 表示两个末尾配对。手算时按这个协议跟踪：a 与 a 相等得到 1；到 abcde 和 ace 的最后，e 与 e 相等，所以由 abcd/ac 的答案加一。

\textbf{代码落点。} 关键是表多开一行一列表示空前缀，循环从 1 开始，访问字符用 text[i-1]。关键代码行必须能逐句翻译成“旧状态怎样变成新状态”，否则就只是模板。

\textbf{正确性。} 任意 LCS 对末尾字符只有配对或至少丢掉一边两类；转移覆盖这两类并取最大。

\textbf{边界实验。} 先用空串、完全相同、无公共字符、子序列不是子串做反例检查。实验时覆盖空串、完全相同、无公共字符、子序列不是子串，先打印完整 DP 表，再比较空间压缩版本。若压缩版不同，优先检查遍历顺序和旧值保存。

\textbf{复杂度变化。} 暴力指数级；DP O(mn) 时间和空间，可滚动到 O(min(m, n)) 空间。

\textbf{迁移追问。} 如果题目改成“若要求还原序列，需要从表右下角反向追踪选择”，先判断原状态是否仍然覆盖新答案集合，再决定是微调边界、改 value 语义，还是换算法。

\subsection{LeetCode 139 单词拆分}

\textbf{题目说明。} LeetCode 139 单词拆分 的题面入口要先还原成具体任务：状态表示前缀能否由字典单词拼出。

\textbf{样例入口。} 拿 s=leetcode、字典 [leet, code] 手算，dp[4] 因为 leet 可拆而为真，dp[8] 又因为 dp[4] 且 code 在字典里为真。

\textbf{暴力基线。} 暴力递归从字符串开头切一个字典单词，剩余后缀继续递归；尝试所有切法。

\textbf{暴力过程。} 以 leetcode 和字典 [leet, code] 为例，前缀 leet 在字典中，剩余 code 也在字典中，所以可拆。

\textbf{重复工作。} 题面模型：状态表示前缀能否由字典单词拼出。暴力版的慢点要落到本题：浪费在同一个后缀会由不同切法反复判断。例如某个下标 i 之后能否拆分，只取决于 i，不取决于前面怎么切。后面的优化只消掉这类重复工作，不能改变答案集合。

\textbf{优化条件。} 本题的关键观察是：枚举最后一个单词的起点，左侧前缀可达且子串在字典中则当前可达。这句话必须能翻译成可维护状态；如果题目条件改变后观察不再成立，就不能继续套原来的写法。

\textbf{相邻状态。} 规律是 dp[i] 表示前 i 个字符能否拆分，转移枚举最后一个单词起点。把这个特例继续拆开看：先问暴力搜索里哪些不同路径到达同一个后续问题，再给这个后续问题命名为状态。状态必须包含未来决策需要的全部信息，转移只从更小且已经算清的状态来。

\textbf{状态复用。} dp[i] 表示 s[0..i) 是否能被字典拆分。若存在 j<i，使 dp[j] 为真且 s[j..i) 在字典中，则 dp[i]=true。手算时只改被新输入影响的那一小部分，而不是回头重算完整候选。

\textbf{指针和字段。} 状态字段要能说清：i 是当前前缀右端，j 是最后一个单词起点，word\_set 用于 O(1) 查询子串是否为单词。手算时按这个协议跟踪：leetcode 中 dp[4] 因 leet 为真；随后 i=8 时 j=4，dp[4] 为真且 code 在字典中，所以 dp[8]=true。

\textbf{代码落点。} 关键是 dp[0]=true 表示空前缀可拆。热循环里 substr 会复制字符串，工程版可按最大单词长度限制 j 或用 Trie。关键代码行必须能逐句翻译成“旧状态怎样变成新状态”，否则就只是模板。

\textbf{正确性。} 任意合法拆分都有最后一个单词 s[j..i)，其左侧前缀必须可拆；转移枚举了所有可能最后单词起点。

\textbf{边界实验。} 先用空串、重复单词、长单词、substr 复制成本做反例检查。实验时覆盖空串、重复单词、长单词、substr 复制成本，先打印完整 DP 表，再比较空间压缩版本。若压缩版不同，优先检查遍历顺序和旧值保存。

\textbf{复杂度变化。} 暴力指数级；DP 有 O(\ensuremath{n^{2}}) 个切分检查，若 substr 计入成本可到 O(\ensuremath{n^{3}})，空间 O(n+字典)。

\textbf{迁移追问。} 如果题目改成“若要输出所有句子，需要 DFS 加记忆化或前驱列表”，先判断原状态是否仍然覆盖新答案集合，再决定是微调边界、改 value 语义，还是换算法。

\subsection{LeetCode 115 不同的子序列}

\textbf{题目说明。} LeetCode 115 不同的子序列 的题面入口要先还原成具体任务：从源串中删除若干字符后形成目标串，本质是两个前缀之间的计数问题。

\textbf{样例入口。} 拿 s=bab、t=ba 手算。扫描到最后一个 b 时，它不能匹配 t 的 a，只能继承不用它的方案；扫描到 a 时，有两类方案：用这个 a 匹配 t 的末尾，或不用它。

\textbf{暴力基线。} 慢版本先写递归搜索树：节点表示处理位置、剩余资源和当前选择。只要不同路径反复到达同一个节点，就说明需要把这个节点命名成 DP 状态。放到本题就是：先按“从源串中删除若干字符后形成目标串，本质是两个前缀之间的计数问题”完整试慢版本；样例里已经能看到“规律是字符相等时计数相加，不等时只继承不用源字符的方案”，所以优化前必须先能说清慢版本枚举了哪些候选。

\textbf{暴力过程。} 拿 s=bab、t=ba 手算。扫描到最后一个 b 时，它不能匹配 t 的 a，只能继承不用它的方案；扫描到 a 时，有两类方案：用这个 a 匹配 t 的末尾，或不用它。

\textbf{重复工作。} 题面模型：从源串中删除若干字符后形成目标串，本质是两个前缀之间的计数问题。暴力版的慢点要落到本题：慢版本会在“从源串中删除若干字符后形成目标串，本质是两个前缀之间的计数问题”的选择树里反复到达相同参数。样例规律是“规律是字符相等时计数相加，不等时只继承不用源字符的方案”，说明这个参数组合可以命名成状态，算过之后后面直接复用。后面的优化只消掉这类重复工作，不能改变答案集合。

\textbf{优化条件。} 本题的关键观察是：状态 dp[i][j] 表示源串前 i 个字符中有多少个子序列等于目标前 j 个字符。这句话必须能翻译成可维护状态；如果题目条件改变后观察不再成立，就不能继续套原来的写法。

\textbf{相邻状态。} 规律是字符相等时计数相加，不等时只继承不用源字符的方案。把这个特例继续拆开看：先问暴力搜索里哪些不同路径到达同一个后续问题，再给这个后续问题命名为状态。状态必须包含未来决策需要的全部信息，转移只从更小且已经算清的状态来。

\textbf{状态复用。} 把重复递归节点命名为状态；遍历顺序只是在保证依赖状态已经算完。本题落点是：规律是字符相等时计数相加，不等时只继承不用源字符的方案。手算时只改被新输入影响的那一小部分，而不是回头重算完整候选。

\textbf{指针和字段。} 状态字段要能说清：状态下标、状态值语义、初始化、转移来源、遍历顺序；空间压缩时还要指出读到的是旧值还是新值。手算时按这个协议跟踪：拿 s=bab、t=ba 手算。扫描到最后一个 b 时，它不能匹配 t 的 a，只能继承不用它的方案；扫描到 a 时，有两类方案：用这个 a 匹配 t 的末尾，或不用它。然后按协议复查：手算时先写一个很小的 DP 表或记忆化搜索记录。每个格子旁边写它来自哪些更小状态。

\textbf{代码落点。} 先写未压缩版本校验转移；压缩前后用小表对拍；不可达哨兵参与加法前要检查溢出。关键代码行必须能逐句翻译成“旧状态怎样变成新状态”，否则就只是模板。

\textbf{正确性。} 证明从最后一步开始：任意最优解的最后一步都在转移枚举中，转移来源已经由更小状态正确计算。

\textbf{边界实验。} 先用空目标串计数为一、源串为空、重复字符导致计数增长、计数可能超过 int做反例检查。实验时覆盖空目标串计数为一、源串为空、重复字符导致计数增长、计数可能超过 int，先打印完整 DP 表，再比较空间压缩版本。若压缩版不同，优先检查遍历顺序和旧值保存。

\textbf{复杂度变化。} 暴力递归会重复计算相同子问题，常见指数级；DP 把每个状态算一次，时间是状态数乘单个转移成本，空间是状态表大小或压缩后的大小。

\textbf{迁移追问。} 如果题目改成“若只问是否存在目标子序列，可以用双指针，不需要二维计数 DP”，先判断原状态是否仍然覆盖新答案集合，再决定是微调边界、改 value 语义，还是换算法。

\subsection{LeetCode 312 戳气球}

\textbf{题目说明。} LeetCode 312 戳气球 的题面入口要先还原成具体任务：先戳哪个会让邻居持续变化，改成最后戳某个气球后区间边界才稳定。

\textbf{样例入口。} 拿 [3, 1, 5] 手算，若先戳 1，左右邻居会变；若改问某个区间里最后戳谁，最后一刻左右边界固定，收益可拆成左右两个子区间。

\textbf{暴力基线。} 慢版本先写递归搜索树：节点表示处理位置、剩余资源和当前选择。只要不同路径反复到达同一个节点，就说明需要把这个节点命名成 DP 状态。放到本题就是：先按“先戳哪个会让邻居持续变化，改成最后戳某个气球后区间边界才稳定”完整试慢版本；样例里已经能看到“规律是区间 DP 枚举最后戳的 mid，而不是第一个戳的气球”，所以优化前必须先能说清慢版本枚举了哪些候选。

\textbf{暴力过程。} 拿 [3, 1, 5] 手算，若先戳 1，左右邻居会变；若改问某个区间里最后戳谁，最后一刻左右边界固定，收益可拆成左右两个子区间。

\textbf{重复工作。} 题面模型：先戳哪个会让邻居持续变化，改成最后戳某个气球后区间边界才稳定。暴力版的慢点要落到本题：慢版本会在“先戳哪个会让邻居持续变化，改成最后戳某个气球后区间边界才稳定”的选择树里反复到达相同参数。样例规律是“规律是区间 DP 枚举最后戳的 mid，而不是第一个戳的气球”，说明这个参数组合可以命名成状态，算过之后后面直接复用。后面的优化只消掉这类重复工作，不能改变答案集合。

\textbf{优化条件。} 本题的关键观察是：区间 DP 定义开区间内全部戳完的最大收益，枚举最后一个被戳的气球作为切分点。这句话必须能翻译成可维护状态；如果题目条件改变后观察不再成立，就不能继续套原来的写法。

\textbf{相邻状态。} 规律是区间 DP 枚举最后戳的 mid，而不是第一个戳的气球。把这个特例继续拆开看：先问暴力搜索里哪些不同路径到达同一个后续问题，再给这个后续问题命名为状态。状态必须包含未来决策需要的全部信息，转移只从更小且已经算清的状态来。

\textbf{状态复用。} 把重复递归节点命名为状态；遍历顺序只是在保证依赖状态已经算完。本题落点是：规律是区间 DP 枚举最后戳的 mid，而不是第一个戳的气球。手算时只改被新输入影响的那一小部分，而不是回头重算完整候选。

\textbf{指针和字段。} 状态字段要能说清：状态下标、状态值语义、初始化、转移来源、遍历顺序；空间压缩时还要指出读到的是旧值还是新值。手算时按这个协议跟踪：拿 [3, 1, 5] 手算，若先戳 1，左右邻居会变；若改问某个区间里最后戳谁，最后一刻左右边界固定，收益可拆成左右两个子区间。然后按协议复查：手算时先写一个很小的 DP 表或记忆化搜索记录。每个格子旁边写它来自哪些更小状态。

\textbf{代码落点。} 先写未压缩版本校验转移；压缩前后用小表对拍；不可达哨兵参与加法前要检查溢出。关键代码行必须能逐句翻译成“旧状态怎样变成新状态”，否则就只是模板。

\textbf{正确性。} 证明从最后一步开始：任意最优解的最后一步都在转移枚举中，转移来源已经由更小状态正确计算。

\textbf{边界实验。} 先用补一的虚拟边界、区间长度遍历顺序、空区间收益为零、乘积可能较大做反例检查。实验时覆盖补一的虚拟边界、区间长度遍历顺序、空区间收益为零、乘积可能较大，先打印完整 DP 表，再比较空间压缩版本。若压缩版不同，优先检查遍历顺序和旧值保存。

\textbf{复杂度变化。} 暴力递归会重复计算相同子问题，常见指数级；DP 把每个状态算一次，时间是状态数乘单个转移成本，空间是状态表大小或压缩后的大小。

\textbf{迁移追问。} 如果题目改成“很多区间 DP 都要换成最后一步思考，避免中间操作改变边界”，先判断原状态是否仍然覆盖新答案集合，再决定是微调边界、改 value 语义，还是换算法。

\subsection{LeetCode 516 最长回文子序列}

\textbf{题目说明。} LeetCode 516 最长回文子序列 的题面入口要先还原成具体任务：子序列允许跳过字符，区间两端是否相等决定能否同时作为回文两端。

\textbf{样例入口。} 拿 s=bbbab 手算，两端 b 相等时可以同时作为回文两端，内部 bba 的最长回文再加 2；若两端不同，就只能丢左端或丢右端取较大。

\textbf{暴力基线。} 慢版本先写递归搜索树：节点表示处理位置、剩余资源和当前选择。只要不同路径反复到达同一个节点，就说明需要把这个节点命名成 DP 状态。放到本题就是：先按“子序列允许跳过字符，区间两端是否相等决定能否同时作为回文两端”完整试慢版本；样例里已经能看到“规律是 dp[left][right] 表示区间最长回文子序列，两端相等用内部区间加二，不等时取两个缩小区间的最大值”，所以优化前必须先能说清慢版本枚举了哪些候选。

\textbf{暴力过程。} 拿 s=bbbab 手算，两端 b 相等时可以同时作为回文两端，内部 bba 的最长回文再加 2；若两端不同，就只能丢左端或丢右端取较大。

\textbf{重复工作。} 题面模型：子序列允许跳过字符，区间两端是否相等决定能否同时作为回文两端。暴力版的慢点要落到本题：慢版本会在“子序列允许跳过字符，区间两端是否相等决定能否同时作为回文两端”的选择树里反复到达相同参数。样例规律是“规律是 dp[left][right] 表示区间最长回文子序列，两端相等用内部区间加二，不等时取两个缩小区间的最大值”，说明这个参数组合可以命名成状态，算过之后后面直接复用。后面的优化只消掉这类重复工作，不能改变答案集合。

\textbf{优化条件。} 本题的关键观察是：dp[left][right] 表示区间最长回文子序列，两端相等取内部加二，否则丢左或丢右取最大。这句话必须能翻译成可维护状态；如果题目条件改变后观察不再成立，就不能继续套原来的写法。

\textbf{相邻状态。} 规律是 dp[left][right] 表示区间最长回文子序列，两端相等用内部区间加二，不等时取两个缩小区间的最大值。把这个特例继续拆开看：先问暴力搜索里哪些不同路径到达同一个后续问题，再给这个后续问题命名为状态。状态必须包含未来决策需要的全部信息，转移只从更小且已经算清的状态来。

\textbf{状态复用。} 把重复递归节点命名为状态；遍历顺序只是在保证依赖状态已经算完。本题落点是：规律是 dp[left][right] 表示区间最长回文子序列，两端相等用内部区间加二，不等时取两个缩小区间的最大值。手算时只改被新输入影响的那一小部分，而不是回头重算完整候选。

\textbf{指针和字段。} 状态字段要能说清：状态下标、状态值语义、初始化、转移来源、遍历顺序；空间压缩时还要指出读到的是旧值还是新值。手算时按这个协议跟踪：拿 s=bbbab 手算，两端 b 相等时可以同时作为回文两端，内部 bba 的最长回文再加 2；若两端不同，就只能丢左端或丢右端取较大。然后按协议复查：手算时先写一个很小的 DP 表或记忆化搜索记录。每个格子旁边写它来自哪些更小状态。

\textbf{代码落点。} 先写未压缩版本校验转移；压缩前后用小表对拍；不可达哨兵参与加法前要检查溢出。关键代码行必须能逐句翻译成“旧状态怎样变成新状态”，否则就只是模板。

\textbf{正确性。} 证明从最后一步开始：任意最优解的最后一步都在转移枚举中，转移来源已经由更小状态正确计算。

\textbf{边界实验。} 先用空串、单字符、两端相等但内部为空、全相同字符、子串和子序列混淆做反例检查。实验时覆盖空串、单字符、两端相等但内部为空、全相同字符、子串和子序列混淆，先打印完整 DP 表，再比较空间压缩版本。若压缩版不同，优先检查遍历顺序和旧值保存。

\textbf{复杂度变化。} 暴力递归会重复计算相同子问题，常见指数级；DP 把每个状态算一次，时间是状态数乘单个转移成本，空间是状态表大小或压缩后的大小。

\textbf{迁移追问。} 如果题目改成“若求最长回文子串，状态必须保持连续性，不能用这个跳过模型”，先判断原状态是否仍然覆盖新答案集合，再决定是微调边界、改 value 语义，还是换算法。

\subsection{LeetCode 847 访问所有节点的最短路径}

\textbf{题目说明。} LeetCode 847 访问所有节点的最短路径 的题面入口要先还原成具体任务：访问集合影响未来能力，同一节点配不同 mask 不是同一状态。

\textbf{样例入口。} 拿 graph=[[1, 2, 3], [0], [0], [0]] 手算，从中心 0 出发访问三个叶子需要来回走，长度为 4；如果只按节点 visited，到 0 时会把访问了不同叶子的状态混掉。

\textbf{暴力基线。} 慢版本先按题意画出完整状态图，用普通搜索跑通可达性。这个过程用于确认访问集合影响未来能力，同一节点配不同 mask 不是同一状态的节点和边，之后再讨论访问标记、层次或代价顺序。放到本题就是：先按“访问集合影响未来能力，同一节点配不同 mask 不是同一状态”完整试慢版本；样例里已经能看到“规律是 BFS 状态必须是 (node, mask)，从所有起点多源出发，第一次到全集 mask 就是最短”，所以优化前必须先能说清慢版本枚举了哪些候选。

\textbf{暴力过程。} 拿 graph=[[1, 2, 3], [0], [0], [0]] 手算，从中心 0 出发访问三个叶子需要来回走，长度为 4；如果只按节点 visited，到 0 时会把访问了不同叶子的状态混掉。

\textbf{重复工作。} 题面模型：访问集合影响未来能力，同一节点配不同 mask 不是同一状态。暴力版的慢点要落到本题：慢版本会围绕“访问集合影响未来能力，同一节点配不同 mask 不是同一状态”反复展开同一批状态。样例规律是“规律是 BFS 状态必须是 (node, mask)，从所有起点多源出发，第一次到全集 mask 就是最短”，说明必须把节点、边、访问标记和资源维度一次建清；同一状态不应重复入队，不同资源状态也不能误合并。后面的优化只消掉这类重复工作，不能改变答案集合。

\textbf{优化条件。} 本题的关键观察是：从所有节点以各自单点 mask 多源入队，BFS 扩展 (node, mask)，第一次到全集 mask 即最短。这句话必须能翻译成可维护状态；如果题目条件改变后观察不再成立，就不能继续套原来的写法。

\textbf{相邻状态。} 规律是 BFS 状态必须是 (node, mask)，从所有起点多源出发，第一次到全集 mask 就是最短。把这个特例继续拆开看：状态节点不一定等于原始位置，还可能包含钥匙、剩余资源、时间或访问集合。visited 只能合并未来能力完全相同的状态，否则会把合法路径剪掉。

\textbf{状态复用。} 把重复搜索压成访问标记、距离数组或拓扑状态；邻居生成也要从全量扫描变成结构化枚举。本题落点是：规律是 BFS 状态必须是 (node, mask)，从所有起点多源出发，第一次到全集 mask 就是最短。手算时只改被新输入影响的那一小部分，而不是回头重算完整候选。

\textbf{指针和字段。} 状态字段要能说清：状态编号、邻接生成方式、访问标记、距离或层数、队列内容；若状态含资源，visited 不能只按位置记录。手算时按这个协议跟踪：拿 graph=[[1, 2, 3], [0], [0], [0]] 手算，从中心 0 出发访问三个叶子需要来回走，长度为 4；如果只按节点 visited，到 0 时会把访问了不同叶子的状态混掉。然后按协议复查：手算时画队列或栈，状态必须包含题目需要的所有资源。入队时标记还是出队时标记，要和去重语义一致。

\textbf{代码落点。} 入队时标记还是出队时标记必须一致；方向数组集中定义；多源 BFS 要一次性把所有源点放入初始队列。关键代码行必须能逐句翻译成“旧状态怎样变成新状态”，否则就只是模板。

\textbf{正确性。} 证明从可达性开始：搜索覆盖所有合法边能到达的状态，访问标记只删除重复状态而不删除不同未来能力。

\textbf{边界实验。} 先用单节点、完全图、星形图、重复经过节点、visited 必须包含 mask做反例检查。实验时覆盖单节点、完全图、星形图、重复经过节点、visited 必须包含 mask，并打印入队时的完整状态。若同一位置但资源不同的状态被合并，很多图题会出现隐蔽错误。

\textbf{复杂度变化。} 暴力重复从多个起点搜索会反复遍历图；正确建模和访问标记后，BFS 或 DFS 通常按节点数加边数计算，状态图题则按完整状态数计算。

\textbf{迁移追问。} 如果题目改成“若边有权，状态图要改用 Dijkstra”，先判断原状态是否仍然覆盖新答案集合，再决定是微调边界、改 value 语义，还是换算法。

\subsection{LeetCode 698 划分为 k 个相等的子集}

\textbf{题目说明。} LeetCode 698 划分为 k 个相等的子集 的题面入口要先还原成具体任务：每个元素要分配到 k 个桶中，所有桶目标和相同。

\textbf{样例入口。} 拿 [4, 3, 2, 3, 5, 2, 1]、k=4 手算，总和 20，每桶目标 5；先放 5 单独成桶，再尝试 4+1、3+2。

\textbf{暴力基线。} 慢版本就是完整搜索树：每层列出选择列表，路径到达终止条件时检查答案。它先保证覆盖所有可能，再把明显无效或重复的分支剪掉。放到本题就是：先按“每个元素要分配到 k 个桶中，所有桶目标和相同”完整试慢版本；样例里已经能看到“规律是回溯状态是桶剩余容量和已用元素，排序后先放大数能更早剪枝”，所以优化前必须先能说清慢版本枚举了哪些候选。

\textbf{暴力过程。} 拿 [4, 3, 2, 3, 5, 2, 1]、k=4 手算，总和 20，每桶目标 5；先放 5 单独成桶，再尝试 4+1、3+2。

\textbf{重复工作。} 题面模型：每个元素要分配到 k 个桶中，所有桶目标和相同。暴力版的慢点要落到本题：慢版本会围绕“每个元素要分配到 k 个桶中，所有桶目标和相同”展开完整搜索树。样例规律是“规律是回溯状态是桶剩余容量和已用元素，排序后先放大数能更早剪枝”，说明一层递归必须对应一个明确决策，剪枝只能删除整棵不可能成功或重复的子树。后面的优化只消掉这类重复工作，不能改变答案集合。

\textbf{优化条件。} 本题的关键观察是：先判断总和可整除并排序降序，再用桶剩余容量回溯放置元素。这句话必须能翻译成可维护状态；如果题目条件改变后观察不再成立，就不能继续套原来的写法。

\textbf{相邻状态。} 规律是回溯状态是桶剩余容量和已用元素，排序后先放大数能更早剪枝。把这个特例继续拆开看：一层递归对应一个明确决策，路径保存已经做出的选择，选择列表保存仍可能扩展的分支。剪枝前必须能说明被剪掉的整棵子树没有合法答案，而不是只因为它看起来不优。

\textbf{状态复用。} 把完整搜索树加上合法性检查、剩余资源剪枝和同层去重；路径状态进入和退出要成对恢复。本题落点是：规律是回溯状态是桶剩余容量和已用元素，排序后先放大数能更早剪枝。手算时只改被新输入影响的那一小部分，而不是回头重算完整候选。

\textbf{指针和字段。} 状态字段要能说清：路径、起点或使用标记、剩余目标、答案集合、剪枝条件；进入递归和退出递归必须成对恢复。手算时按这个协议跟踪：拿 [4, 3, 2, 3, 5, 2, 1]、k=4 手算，总和 20，每桶目标 5；先放 5 单独成桶，再尝试 4+1、3+2。然后按协议复查：手算时给每层标出选择列表和路径。剪枝发生前要指出被剪分支为什么已经不可能得到合法答案。

\textbf{代码落点。} 剪枝条件写在扩展前；同层去重不要误写成全局去重；保存答案时复制当前路径。关键代码行必须能逐句翻译成“旧状态怎样变成新状态”，否则就只是模板。

\textbf{正确性。} 证明从搜索树覆盖开始：每个合法答案有且只有一条路径，剪枝只删掉不可能成功的分支。

\textbf{边界实验。} 先用存在元素大于目标、重复元素、空桶对称、k 为一做反例检查。实验时覆盖存在元素大于目标、重复元素、空桶对称、k 为一，统计搜索节点数量和剪枝次数。重复元素题要打印同一层跳过哪些值，确认没有误删合法路径。

\textbf{复杂度变化。} 暴力搜索树的规模由输出或选择空间决定；剪枝不改变最坏指数级本质，但能删除不可能成功或重复的分支，复杂度要说明输出规模。

\textbf{迁移追问。} 如果题目改成“状态压缩 DP 也能做，适合用掩码表示已选元素集合”，先判断原状态是否仍然覆盖新答案集合，再决定是微调边界、改 value 语义，还是换算法。

\subsection{LeetCode 864 获取所有钥匙的最短路径}

\textbf{题目说明。} LeetCode 864 获取所有钥匙的最短路径 的题面入口要先还原成具体任务：位置相同但钥匙集合不同，未来能开的门不同。

\textbf{样例入口。} 拿网格 @.a / \#.\# / b.A 手算，走到 a 后状态不是位置变了而是钥匙集合多了一把；同一位置带不同钥匙未来能力不同。

\textbf{暴力基线。} 慢版本先按题意画出完整状态图，用普通搜索跑通可达性。这个过程用于确认位置相同但钥匙集合不同，未来能开的门不同的节点和边，之后再讨论访问标记、层次或代价顺序。放到本题就是：先按“位置相同但钥匙集合不同，未来能开的门不同”完整试慢版本；样例里已经能看到“规律是 BFS 的 visited 必须按 (row, col, keymask) 记录，拿齐所有钥匙第一次出队就是最短步数”，所以优化前必须先能说清慢版本枚举了哪些候选。

\textbf{暴力过程。} 拿网格 @.a / \#.\# / b.A 手算，走到 a 后状态不是位置变了而是钥匙集合多了一把；同一位置带不同钥匙未来能力不同。

\textbf{重复工作。} 题面模型：位置相同但钥匙集合不同，未来能开的门不同。暴力版的慢点要落到本题：慢版本会围绕“位置相同但钥匙集合不同，未来能开的门不同”反复展开同一批状态。样例规律是“规律是 BFS 的 visited 必须按 (row, col, keymask) 记录，拿齐所有钥匙第一次出队就是最短步数”，说明必须把节点、边、访问标记和资源维度一次建清；同一状态不应重复入队，不同资源状态也不能误合并。后面的优化只消掉这类重复工作，不能改变答案集合。

\textbf{优化条件。} 本题的关键观察是：BFS 状态包含行、列和钥匙掩码，访问标记也必须包含掩码。这句话必须能翻译成可维护状态；如果题目条件改变后观察不再成立，就不能继续套原来的写法。

\textbf{相邻状态。} 规律是 BFS 的 visited 必须按 (row, col, keymask) 记录，拿齐所有钥匙第一次出队就是最短步数。把这个特例继续拆开看：状态节点不一定等于原始位置，还可能包含钥匙、剩余资源、时间或访问集合。visited 只能合并未来能力完全相同的状态，否则会把合法路径剪掉。

\textbf{状态复用。} 把重复搜索压成访问标记、距离数组或拓扑状态；邻居生成也要从全量扫描变成结构化枚举。本题落点是：规律是 BFS 的 visited 必须按 (row, col, keymask) 记录，拿齐所有钥匙第一次出队就是最短步数。手算时只改被新输入影响的那一小部分，而不是回头重算完整候选。

\textbf{指针和字段。} 状态字段要能说清：状态编号、邻接生成方式、访问标记、距离或层数、队列内容；若状态含资源，visited 不能只按位置记录。手算时按这个协议跟踪：拿网格 @.a / \#.\# / b.A 手算，走到 a 后状态不是位置变了而是钥匙集合多了一把；同一位置带不同钥匙未来能力不同。然后按协议复查：手算时画队列或栈，状态必须包含题目需要的所有资源。入队时标记还是出队时标记，要和去重语义一致。

\textbf{代码落点。} 入队时标记还是出队时标记必须一致；方向数组集中定义；多源 BFS 要一次性把所有源点放入初始队列。关键代码行必须能逐句翻译成“旧状态怎样变成新状态”，否则就只是模板。

\textbf{正确性。} 证明从可达性开始：搜索覆盖所有合法边能到达的状态，访问标记只删除重复状态而不删除不同未来能力。

\textbf{边界实验。} 先用起点旁边有门、钥匙顺序、不可达、钥匙数量决定掩码大小做反例检查。实验时覆盖起点旁边有门、钥匙顺序、不可达、钥匙数量决定掩码大小，并打印入队时的完整状态。若同一位置但资源不同的状态被合并，很多图题会出现隐蔽错误。

\textbf{复杂度变化。} 暴力重复从多个起点搜索会反复遍历图；正确建模和访问标记后，BFS 或 DFS 通常按节点数加边数计算，状态图题则按完整状态数计算。

\textbf{迁移追问。} 如果题目改成“这是状态图维度不能只按位置去重的典型题”，先判断原状态是否仍然覆盖新答案集合，再决定是微调边界、改 value 语义，还是换算法。

\subsection{LeetCode 1425 带限制的子序列和}

\textbf{题目说明。} LeetCode 1425 带限制的子序列和 的题面入口要先还原成具体任务：dp[i] 只依赖前 k 个位置的最大 dp 值，转移范围是滑动窗口。

\textbf{样例入口。} 拿 nums=[10, 2, -10, 5, 20]、k=2 手算，dp[3] 的前驱只能在下标 1 和 2 中选，dp[4] 又能接 dp[3] 得到 37。暴力每个 i 回看前 k 个 dp；单调队列保留窗口内 dp 最大的下标。

\textbf{暴力基线。} 慢版本先写递归搜索树：节点表示处理位置、剩余资源和当前选择。只要不同路径反复到达同一个节点，就说明需要把这个节点命名成 DP 状态。放到本题就是：先按“dp[i] 只依赖前 k 个位置的最大 dp 值，转移范围是滑动窗口”完整试慢版本；样例里已经能看到“规律是转移 dp[i]=nums[i]+max(0, max dp[j])，其中 j 在 i-k 到 i-1 内”，所以优化前必须先能说清慢版本枚举了哪些候选。

\textbf{暴力过程。} 拿 nums=[10, 2, -10, 5, 20]、k=2 手算，dp[3] 的前驱只能在下标 1 和 2 中选，dp[4] 又能接 dp[3] 得到 37。暴力每个 i 回看前 k 个 dp；单调队列保留窗口内 dp 最大的下标。

\textbf{重复工作。} 题面模型：dp[i] 只依赖前 k 个位置的最大 dp 值，转移范围是滑动窗口。暴力版的慢点要落到本题：慢版本会在“dp[i] 只依赖前 k 个位置的最大 dp 值，转移范围是滑动窗口”的选择树里反复到达相同参数。样例规律是“规律是转移 dp[i]=nums[i]+max(0, max dp[j])，其中 j 在 i-k 到 i-1 内”，说明这个参数组合可以命名成状态，算过之后后面直接复用。后面的优化只消掉这类重复工作，不能改变答案集合。

\textbf{优化条件。} 本题的关键观察是：单调队列保存候选 dp 下标，过期弹队首，被更大 dp 支配的候选从队尾删除。这句话必须能翻译成可维护状态；如果题目条件改变后观察不再成立，就不能继续套原来的写法。

\textbf{相邻状态。} 规律是转移 dp[i]=nums[i]+max(0, max dp[j])，其中 j 在 i-k 到 i-1 内。把这个特例继续拆开看：先问暴力搜索里哪些不同路径到达同一个后续问题，再给这个后续问题命名为状态。状态必须包含未来决策需要的全部信息，转移只从更小且已经算清的状态来。

\textbf{状态复用。} 把重复递归节点命名为状态；遍历顺序只是在保证依赖状态已经算完。本题落点是：规律是转移 dp[i]=nums[i]+max(0, max dp[j])，其中 j 在 i-k 到 i-1 内。手算时只改被新输入影响的那一小部分，而不是回头重算完整候选。

\textbf{指针和字段。} 状态字段要能说清：状态下标、状态值语义、初始化、转移来源、遍历顺序；空间压缩时还要指出读到的是旧值还是新值。手算时按这个协议跟踪：拿 nums=[10, 2, -10, 5, 20]、k=2 手算，dp[3] 的前驱只能在下标 1 和 2 中选，dp[4] 又能接 dp[3] 得到 37。暴力每个 i 回看前 k 个 dp；单调队列保留窗口内 dp 最大的下标。然后按协议复查：手算时先写一个很小的 DP 表或记忆化搜索记录。每个格子旁边写它来自哪些更小状态。

\textbf{代码落点。} 先写未压缩版本校验转移；压缩前后用小表对拍；不可达哨兵参与加法前要检查溢出。关键代码行必须能逐句翻译成“旧状态怎样变成新状态”，否则就只是模板。

\textbf{正确性。} 证明从最后一步开始：任意最优解的最后一步都在转移枚举中，转移来源已经由更小状态正确计算。

\textbf{边界实验。} 先用全负数、k 为一、k 大于数组长度、答案必须非空、队列过期顺序做反例检查。实验时覆盖全负数、k 为一、k 大于数组长度、答案必须非空、队列过期顺序，先打印完整 DP 表，再比较空间压缩版本。若压缩版不同，优先检查遍历顺序和旧值保存。

\textbf{复杂度变化。} 暴力递归会重复计算相同子问题，常见指数级；DP 把每个状态算一次，时间是状态数乘单个转移成本，空间是状态表大小或压缩后的大小。

\textbf{迁移追问。} 如果题目改成“若窗口内需要最小值或其他可合并统计，队列维护规则要重证”，先判断原状态是否仍然覆盖新答案集合，再决定是微调边界、改 value 语义，还是换算法。

\subsection{LeetCode 337 打家劫舍 III}

\textbf{题目说明。} LeetCode 337 打家劫舍 III 的题面入口要先还原成具体任务：树上相邻节点不能同时选，每个节点需要偷和不偷两个状态。

\textbf{样例入口。} 拿树 [3, 2, 3, null, 3, null, 1] 手算，偷根 3 时不能偷两个孩子，只能接孙子；不偷根时可以分别取左右孩子的最优。

\textbf{暴力基线。} 慢版本在每个节点重新扫描子树或路径，先求出局部答案。重复扫描暴露了真正状态：每个节点只需要从孩子或父亲那里拿到少量信息。放到本题就是：先按“树上相邻节点不能同时选，每个节点需要偷和不偷两个状态”完整试慢版本；样例里已经能看到“规律是每个节点返回两个状态：偷当前和不偷当前，父节点再组合”，所以优化前必须先能说清慢版本枚举了哪些候选。

\textbf{暴力过程。} 拿树 [3, 2, 3, null, 3, null, 1] 手算，偷根 3 时不能偷两个孩子，只能接孙子；不偷根时可以分别取左右孩子的最优。

\textbf{重复工作。} 题面模型：树上相邻节点不能同时选，每个节点需要偷和不偷两个状态。暴力版的慢点要落到本题：慢版本会在树上重复确认“树上相邻节点不能同时选，每个节点需要偷和不偷两个状态”。样例规律是“规律是每个节点返回两个状态：偷当前和不偷当前，父节点再组合”，说明父子之间真正需要传递的是高度、上下界、命中状态、层号或两类选择值，而不是反复扫描整棵子树。后面的优化只消掉这类重复工作，不能改变答案集合。

\textbf{优化条件。} 本题的关键观察是：后序汇总：偷当前则孩子不能偷，不偷当前则孩子可偷可不偷。这句话必须能翻译成可维护状态；如果题目条件改变后观察不再成立，就不能继续套原来的写法。

\textbf{相邻状态。} 规律是每个节点返回两个状态：偷当前和不偷当前，父节点再组合。把这个特例继续拆开看：节点向父亲返回的量和全局答案通常不是同一个量。先写叶子或空节点的返回值，再写父节点如何合并左右孩子，递归式才有边界基础。

\textbf{状态复用。} 把对子树的重复扫描压成返回值或队列状态；每个节点只合并孩子给出的稳定信息。本题落点是：规律是每个节点返回两个状态：偷当前和不偷当前，父节点再组合。手算时只改被新输入影响的那一小部分，而不是回头重算完整候选。

\textbf{指针和字段。} 状态字段要能说清：节点指针、传入约束或返回值、当前答案、层号或路径状态；空节点的返回值必须和状态定义匹配。手算时按这个协议跟踪：拿树 [3, 2, 3, null, 3, null, 1] 手算，偷根 3 时不能偷两个孩子，只能接孙子；不偷根时可以分别取左右孩子的最优。然后按协议复查：手算时从叶子开始写返回值，或者从根开始写传入约束。不要只写访问顺序，要写每条边上传递的信息。

\textbf{代码落点。} 递归版本先处理空节点；迭代版本的栈元素要带完整状态；全负值、重复值和极端深树要单独测。关键代码行必须能逐句翻译成“旧状态怎样变成新状态”，否则就只是模板。

\textbf{正确性。} 证明从子树归纳或层序不变量开始：孩子状态正确时父节点合并正确；若用队列，当前轮队列必须正好保存当前层节点。

\textbf{边界实验。} 先用空树、负值是否可能、链式树、递归深度做反例检查。实验时覆盖空树、负值是否可能、链式树、递归深度，尤其关注空节点、单节点、链式树和全负值。递归版本和显式栈版本可以互相对拍。

\textbf{复杂度变化。} 暴力在每个节点重复扫描子树或反复寻找层边界可能退化到平方级；后序汇总、显式栈或层序遍历让每个节点处理常数次，通常是线性时间。

\textbf{迁移追问。} 如果题目改成“可用显式栈后序保存节点状态，避免调用栈”，先判断原状态是否仍然覆盖新答案集合，再决定是微调边界、改 value 语义，还是换算法。

\subsection{LeetCode 834 树中距离之和}

\textbf{题目说明。} LeetCode 834 树中距离之和 的题面入口要先还原成具体任务：每个节点作为根的距离和可以由相邻换根关系转移。

\textbf{样例入口。} 拿 n=6 的样例树手算，以 0 为根先算子树大小和距离和；把根从 0 换到 2 时，2 子树内 4 个点距离都减一，外部 2 个点距离都加一。

\textbf{暴力基线。} 慢版本在每个节点重新扫描子树或路径，先求出局部答案。重复扫描暴露了真正状态：每个节点只需要从孩子或父亲那里拿到少量信息。放到本题就是：先按“每个节点作为根的距离和可以由相邻换根关系转移”完整试慢版本；样例里已经能看到“规律是先求一个根的答案，再沿边按“子树内减少、子树外增加”的关系换根”，所以优化前必须先能说清慢版本枚举了哪些候选。

\textbf{暴力过程。} 拿 n=6 的样例树手算，以 0 为根先算子树大小和距离和；把根从 0 换到 2 时，2 子树内 4 个点距离都减一，外部 2 个点距离都加一。

\textbf{重复工作。} 题面模型：每个节点作为根的距离和可以由相邻换根关系转移。暴力版的慢点要落到本题：慢版本会在树上重复确认“每个节点作为根的距离和可以由相邻换根关系转移”。样例规律是“规律是先求一个根的答案，再沿边按“子树内减少、子树外增加”的关系换根”，说明父子之间真正需要传递的是高度、上下界、命中状态、层号或两类选择值，而不是反复扫描整棵子树。后面的优化只消掉这类重复工作，不能改变答案集合。

\textbf{优化条件。} 本题的关键观察是：先后序计算子树大小和根节点答案，再从父到子换根：子树内距离减一，子树外距离加一。这句话必须能翻译成可维护状态；如果题目条件改变后观察不再成立，就不能继续套原来的写法。

\textbf{相邻状态。} 规律是先求一个根的答案，再沿边按“子树内减少、子树外增加”的关系换根。把这个特例继续拆开看：节点向父亲返回的量和全局答案通常不是同一个量。先写叶子或空节点的返回值，再写父节点如何合并左右孩子，递归式才有边界基础。

\textbf{状态复用。} 把对子树的重复扫描压成返回值或队列状态；每个节点只合并孩子给出的稳定信息。本题落点是：规律是先求一个根的答案，再沿边按“子树内减少、子树外增加”的关系换根。手算时只改被新输入影响的那一小部分，而不是回头重算完整候选。

\textbf{指针和字段。} 状态字段要能说清：节点指针、传入约束或返回值、当前答案、层号或路径状态；空节点的返回值必须和状态定义匹配。手算时按这个协议跟踪：拿 n=6 的样例树手算，以 0 为根先算子树大小和距离和；把根从 0 换到 2 时，2 子树内 4 个点距离都减一，外部 2 个点距离都加一。然后按协议复查：手算时从叶子开始写返回值，或者从根开始写传入约束。不要只写访问顺序，要写每条边上传递的信息。

\textbf{代码落点。} 递归版本先处理空节点；迭代版本的栈元素要带完整状态；全负值、重复值和极端深树要单独测。关键代码行必须能逐句翻译成“旧状态怎样变成新状态”，否则就只是模板。

\textbf{正确性。} 证明从子树归纳或层序不变量开始：孩子状态正确时父节点合并正确；若用队列，当前轮队列必须正好保存当前层节点。

\textbf{边界实验。} 先用单节点、链式树、星形树、n 很大递归栈、边是无向需建双向邻接做反例检查。实验时覆盖单节点、链式树、星形树、n 很大递归栈、边是无向需建双向邻接，尤其关注空节点、单节点、链式树和全负值。递归版本和显式栈版本可以互相对拍。

\textbf{复杂度变化。} 暴力在每个节点重复扫描子树或反复寻找层边界可能退化到平方级；后序汇总、显式栈或层序遍历让每个节点处理常数次，通常是线性时间。

\textbf{迁移追问。} 如果题目改成“若边有权，换根公式要把 size 项乘以边权”，先判断原状态是否仍然覆盖新答案集合，再决定是微调边界、改 value 语义，还是换算法。

\subsection{LeetCode 28 找出字符串中第一个匹配项的下标}

\textbf{题目说明。} LeetCode 28 找出字符串中第一个匹配项的下标 的题面入口要先还原成具体任务：字符串匹配的暴力回退会反复比较已经匹配过的前缀。

\textbf{样例入口。} 拿 haystack=sadbutsad、needle=sad 手算，暴力从每个起点尝试匹配，起点 0 成功；若文本和模式有大量重复前缀，失配后从下一个起点重比会浪费。

\textbf{暴力基线。} 暴力从 haystack 的每个起点尝试匹配 needle，遇到失配就把起点加一，needle 从头再比。

\textbf{暴力过程。} 以 haystack=sadbutsad、needle=sad 为例，起点 0 成功。若换成很多重复前缀的文本，如 aaaaaab 匹配 aaaab，暴力会反复比较相同的 a。

\textbf{重复工作。} 题面模型：字符串匹配的暴力回退会反复比较已经匹配过的前缀。暴力版的慢点要落到本题：浪费在失配后完全丢掉已经匹配的前缀信息。模式串内部若有相同前后缀，下一次匹配可以复用一部分。后面的优化只消掉这类重复工作，不能改变答案集合。

\textbf{优化条件。} 本题的关键观察是：KMP 用前缀函数记录模式串最长相等前后缀，失配时把模式串跳到可复用前缀位置。这句话必须能翻译成可维护状态；如果题目条件改变后观察不再成立，就不能继续套原来的写法。

\textbf{相邻状态。} 规律是 KMP 用模式串前缀函数记录可复用的最长相等前后缀，失配时移动模式而不是回退文本。把这个特例继续拆开看：先标出已经确认的前缀或后缀，再标出未知区；能被丢掉的候选，必须是以后再也不会改变已确认区域语义的候选。这样才算从例子推出数组不变量，而不是只记一次操作。

\textbf{状态复用。} KMP 预处理 lps[i]，表示 needle[0..i] 的最长相等真前后缀长度。匹配时文本指针不回退，模式指针按 lps 跳转。手算时只改被新输入影响的那一小部分，而不是回头重算完整候选。

\textbf{指针和字段。} 状态字段要能说清：i 是文本位置，j 是模式位置，lps[j-1] 是失配后可复用的模式前缀长度。手算时按这个协议跟踪：模式 aaaab 在第 5 个字符失配时，前面 aaa 已经可作为下一轮前缀，不必让文本回到旧起点重新比较。

\textbf{代码落点。} 关键代码是失配时 if (j>0) j=lps[j-1]; else ++i。匹配到 j==m 时返回 i-m。关键代码行必须能逐句翻译成“旧状态怎样变成新状态”，否则就只是模板。

\textbf{正确性。} lps 保存的是已经匹配部分中最长可复用前缀，跳转后所有被跳过的起点都无法比这个前缀匹配更多字符。

\textbf{边界实验。} 先用空模式串、模式比文本长、重复字符、完全不匹配、重叠匹配做反例检查。实验时把空模式串、模式比文本长、重复字符、完全不匹配、重叠匹配列成表，再打印关键下标和有效区间。若某个元素被写入后又被未来步骤破坏，说明区间不变量没有成立。

\textbf{复杂度变化。} 暴力最坏 O(nm)，KMP 预处理 O(m)、匹配 O(n)，空间 O(m)。

\textbf{迁移追问。} 如果题目改成“若只做少量短串匹配，暴力可接受；大量匹配再考虑 KMP、哈希或自动机”，先判断原状态是否仍然覆盖新答案集合，再决定是微调边界、改 value 语义，还是换算法。

\subsection{LeetCode 647 回文子串}

\textbf{题目说明。} LeetCode 647 回文子串 的题面入口要先还原成具体任务：每个回文都有唯一中心或中心间隙。

\textbf{样例入口。} 拿 aaa 手算，以每个字符为中心有 3 个单字符回文，以两个字符中间为中心有 2 个 aa，以中间 a 为中心还有 aaa。

\textbf{暴力基线。} 暴力枚举所有子串，再判断每个子串是否回文，回文就计数。

\textbf{暴力过程。} 以 aaa 为例，暴力会检查 a、aa、aaa、a、aa、a。很多判断都重复比较相同中心附近的字符。

\textbf{重复工作。} 题面模型：每个回文都有唯一中心或中心间隙。暴力版的慢点要落到本题：浪费在每个子串单独判断。每个回文都有唯一中心或中心间隙，从中心扩展能一次数出围绕它的所有回文。后面的优化只消掉这类重复工作，不能改变答案集合。

\textbf{优化条件。} 本题的关键观察是：对每个中心向两边扩展，扩展成功一次就计数一次。这句话必须能翻译成可维护状态；如果题目条件改变后观察不再成立，就不能继续套原来的写法。

\textbf{相邻状态。} 规律是枚举字符中心和间隙中心，每次扩展成功就新增一个回文子串。把这个特例继续拆开看：先标出已经确认的前缀或后缀，再标出未知区；能被丢掉的候选，必须是以后再也不会改变已确认区域语义的候选。这样才算从例子推出数组不变量，而不是只记一次操作。

\textbf{状态复用。} 枚举 2n-1 个中心，每个中心向左右扩展；每扩展成功一次，说明找到一个新回文子串，count++。手算时只改被新输入影响的那一小部分，而不是回头重算完整候选。

\textbf{指针和字段。} 状态字段要能说清：left/right 是当前中心扩展边界，count 是已经找到的回文子串数。手算时按这个协议跟踪：aaa 中以中间 a 为中心可数出 a、aaa；以两个 a 的间隙为中心可数出 aa。

\textbf{代码落点。} 关键是对每个 i 调用 expand(i, i) 和 expand(i, i+1)，expand 每次字符相等就 ++count 再继续外扩。关键代码行必须能逐句翻译成“旧状态怎样变成新状态”，否则就只是模板。

\textbf{正确性。} 每个回文子串对应唯一中心；扩展时每个合法半径都会被计数一次，因此不重不漏。

\textbf{边界实验。} 先用空串、全相同字符、偶数回文、奇数回文做反例检查。实验时把空串、全相同字符、偶数回文、奇数回文列成表，再打印关键下标和有效区间。若某个元素被写入后又被未来步骤破坏，说明区间不变量没有成立。

\textbf{复杂度变化。} 暴力 O(\ensuremath{n^{3}})，中心扩展 O(\ensuremath{n^{2}}) 时间、O(1) 空间。

\textbf{迁移追问。} 如果题目改成“Manacher 算法可线性解决，但中心扩展更适合面试基础”，先判断原状态是否仍然覆盖新答案集合，再决定是微调边界、改 value 语义，还是换算法。

\subsection{LeetCode 5 最长回文子串}

\textbf{题目说明。} LeetCode 5 最长回文子串 的题面入口要先还原成具体任务：回文围绕中心对称，中心可以是字符也可以是两个字符之间。

\textbf{样例入口。} 拿 aba 和 abba 手算，回文不是任意子串都要重查，而是围绕中心向两边同步比较。aba 的中心是 b，abba 的中心在两个 b 中间；每扩一层只新增左右两个字符。

\textbf{暴力基线。} 暴力枚举所有子串 s[left..right]，再用双指针检查它是否回文，记录最长回文子串。

\textbf{暴力过程。} 以 s=babad 为例，暴力会检查 b、ba、bab、baba、babad，再从 a 开始检查 a、ab、aba。bab 和 aba 的内部字符被反复比较。

\textbf{重复工作。} 题面模型：回文围绕中心对称，中心可以是字符也可以是两个字符之间。暴力版的慢点要落到本题：浪费在每个子串都从头判断回文，明明回文的合法性只取决于左右边界是否相等，以及内部是否已经是回文。后面的优化只消掉这类重复工作，不能改变答案集合。

\textbf{优化条件。} 本题的关键观察是：中心扩展避免枚举子串后再逐个检查；每次扩展只比较新增左右边界。这句话必须能翻译成可维护状态；如果题目条件改变后观察不再成立，就不能继续套原来的写法。

\textbf{相邻状态。} 规律是枚举 2n-1 个中心，每个中心只在新增边界相等时继续扩展。把这个特例继续拆开看：先标出已经确认的前缀或后缀，再标出未知区；能被丢掉的候选，必须是以后再也不会改变已确认区域语义的候选。这样才算从例子推出数组不变量，而不是只记一次操作。

\textbf{状态复用。} 把子串枚举改成中心枚举。中心可以是一个字符，也可以是两个字符之间；每次只向左右各扩一格，成功就更新答案边界。手算时只改被新输入影响的那一小部分，而不是回头重算完整候选。

\textbf{指针和字段。} 状态字段要能说清：center 表示回文中心，left/right 是当前扩展边界，best\_left/best\_len 记录最长回文的位置和长度。手算时按这个协议跟踪：aba 从中心 b 开始，先看 b，再扩到 a..a 成功；abba 从两个 b 中间开始，先看 bb，再扩到 a..a 成功。

\textbf{代码落点。} 关键代码是对每个 i 调用 expand(i, i) 和 expand(i, i+1)。expand 里 while left>=0 \&\& right<n \&\& s[left]==s[right]，再更新最长长度。关键代码行必须能逐句翻译成“旧状态怎样变成新状态”，否则就只是模板。

\textbf{正确性。} 任意回文都有唯一中心或中心间隙；从这个中心向外扩展时会经过它的全部边界，所以不会漏掉任何回文子串。

\textbf{边界实验。} 先用偶数长度回文、单字符、全相同字符、多个同长答案做反例检查。实验时把偶数长度回文、单字符、全相同字符、多个同长答案列成表，再打印关键下标和有效区间。若某个元素被写入后又被未来步骤破坏，说明区间不变量没有成立。

\textbf{复杂度变化。} 暴力枚举加检查可到 O(\ensuremath{n^{3}})，中心扩展有 O(n) 个中心、每个中心最多扩 O(n)，时间 O(\ensuremath{n^{2}})，额外空间 O(1)。

\textbf{迁移追问。} 如果题目改成“若要求回文子序列，应转成区间 DP，不能用中心扩展”，先判断原状态是否仍然覆盖新答案集合，再决定是微调边界、改 value 语义，还是换算法。

\subsection{LeetCode 212 单词搜索 II}

\textbf{题目说明。} LeetCode 212 单词搜索 II 的题面入口要先还原成具体任务：多个单词在同一棋盘上搜索，公共前缀可以共享。

\textbf{样例入口。} 拿 board 上有 oath、pea、eat 手算，单词共享前缀时不应对每个词重复全图搜索；从 o 出发沿 Trie 前缀 o-a-t-h 找到 oath。

\textbf{暴力基线。} 慢版本就是完整搜索树：每层列出选择列表，路径到达终止条件时检查答案。它先保证覆盖所有可能，再把明显无效或重复的分支剪掉。放到本题就是：先按“多个单词在同一棋盘上搜索，公共前缀可以共享”完整试慢版本；样例里已经能看到“规律是 Trie 合并所有单词前缀，DFS 走棋盘同时走 Trie，命中单词后记录并可去重”，所以优化前必须先能说清慢版本枚举了哪些候选。

\textbf{暴力过程。} 拿 board 上有 oath、pea、eat 手算，单词共享前缀时不应对每个词重复全图搜索；从 o 出发沿 Trie 前缀 o-a-t-h 找到 oath。

\textbf{重复工作。} 题面模型：多个单词在同一棋盘上搜索，公共前缀可以共享。暴力版的慢点要落到本题：慢版本会围绕“多个单词在同一棋盘上搜索，公共前缀可以共享”展开完整搜索树。样例规律是“规律是 Trie 合并所有单词前缀，DFS 走棋盘同时走 Trie，命中单词后记录并可去重”，说明一层递归必须对应一个明确决策，剪枝只能删除整棵不可能成功或重复的子树。后面的优化只消掉这类重复工作，不能改变答案集合。

\textbf{优化条件。} 本题的关键观察是：先建 Trie，再从每个格子回溯；路径到达单词结束节点就收集答案。这句话必须能翻译成可维护状态；如果题目条件改变后观察不再成立，就不能继续套原来的写法。

\textbf{相邻状态。} 规律是 Trie 合并所有单词前缀，DFS 走棋盘同时走 Trie，命中单词后记录并可去重。把这个特例继续拆开看：一层递归对应一个明确决策，路径保存已经做出的选择，选择列表保存仍可能扩展的分支。剪枝前必须能说明被剪掉的整棵子树没有合法答案，而不是只因为它看起来不优。

\textbf{状态复用。} 把完整搜索树加上合法性检查、剩余资源剪枝和同层去重；路径状态进入和退出要成对恢复。本题落点是：规律是 Trie 合并所有单词前缀，DFS 走棋盘同时走 Trie，命中单词后记录并可去重。手算时只改被新输入影响的那一小部分，而不是回头重算完整候选。

\textbf{指针和字段。} 状态字段要能说清：路径、起点或使用标记、剩余目标、答案集合、剪枝条件；进入递归和退出递归必须成对恢复。手算时按这个协议跟踪：拿 board 上有 oath、pea、eat 手算，单词共享前缀时不应对每个词重复全图搜索；从 o 出发沿 Trie 前缀 o-a-t-h 找到 oath。然后按协议复查：手算时给每层标出选择列表和路径。剪枝发生前要指出被剪分支为什么已经不可能得到合法答案。

\textbf{代码落点。} 剪枝条件写在扩展前；同层去重不要误写成全局去重；保存答案时复制当前路径。关键代码行必须能逐句翻译成“旧状态怎样变成新状态”，否则就只是模板。

\textbf{正确性。} 证明从搜索树覆盖开始：每个合法答案有且只有一条路径，剪枝只删掉不可能成功的分支。

\textbf{边界实验。} 先用重复单词、同一格不能复用、前缀是完整单词、剪枝删除空子树做反例检查。实验时覆盖重复单词、同一格不能复用、前缀是完整单词、剪枝删除空子树，统计搜索节点数量和剪枝次数。重复元素题要打印同一层跳过哪些值，确认没有误删合法路径。

\textbf{复杂度变化。} 暴力搜索树的规模由输出或选择空间决定；剪枝不改变最坏指数级本质，但能删除不可能成功或重复的分支，复杂度要说明输出规模。

\textbf{迁移追问。} 如果题目改成“这是单词搜索从单目标扩展到多目标后的结构升级”，先判断原状态是否仍然覆盖新答案集合，再决定是微调边界、改 value 语义，还是换算法。

\subsection{LeetCode 1044 最长重复子串}

\textbf{题目说明。} LeetCode 1044 最长重复子串 的题面入口要先还原成具体任务：重复子串长度具有可行性单调：存在长度 L 的重复，则更短长度也存在。

\textbf{样例入口。} 拿 banana 手算，重复子串 ana 出现两次，长度 3 可行；若长度 4 不可行，则更长也不可能。

\textbf{暴力基线。} 慢版本按顺序扫描下标、逐个试答案值，或直接检查候选直到遇到目标边界。它先保证候选空间覆盖完整，但每次只排除一个候选，浪费了题目给出的顺序、比较或可行性边界。放到本题就是：先按“重复子串长度具有可行性单调：存在长度 L 的重复，则更短长度也存在”完整试慢版本；样例里已经能看到“规律是二分重复子串长度，用滚动哈希检查同长度子串是否重复，哈希相等时还要防冲突”，所以优化前必须先能说清慢版本枚举了哪些候选。

\textbf{暴力过程。} 拿 banana 手算，重复子串 ana 出现两次，长度 3 可行；若长度 4 不可行，则更长也不可能。

\textbf{重复工作。} 题面模型：重复子串长度具有可行性单调：存在长度 L 的重复，则更短长度也存在。暴力版的慢点要落到本题：慢版本会按顺序试“重复子串长度具有可行性单调：存在长度 L 的重复，则更短长度也存在”里的候选；而样例规律已经指出“规律是二分重复子串长度，用滚动哈希检查同长度子串是否重复，哈希相等时还要防冲突”。一次比较或一次可行性检查能丢掉一整段候选，继续线性试就是浪费。后面的优化只消掉这类重复工作，不能改变答案集合。

\textbf{优化条件。} 本题的关键观察是：二分长度，用滚动哈希检查是否存在两个相同长度子串哈希相同，并在冲突时确认字符串。这句话必须能翻译成可维护状态；如果题目条件改变后观察不再成立，就不能继续套原来的写法。

\textbf{相邻状态。} 规律是二分重复子串长度，用滚动哈希检查同长度子串是否重复，哈希相等时还要防冲突。把这个特例继续拆开看：每次比较 mid 之后，必须能指出哪一侧整体不可能包含目标，或哪一侧整体已经满足可行性。二分的核心不是 mid 公式，而是被丢弃区间确实不含答案。

\textbf{状态复用。} 把逐个试候选改成维护答案所在区间；边界谓词、可行性检查或局部比较给出分支方向，每个分支都必须能排除一侧候选。本题落点是：规律是二分重复子串长度，用滚动哈希检查同长度子串是否重复，哈希相等时还要防冲突。手算时只改被新输入影响的那一小部分，而不是回头重算完整候选。

\textbf{指针和字段。} 状态字段要能说清：left、right、mid、mid 处比较值或检查返回值、答案区间含义、返回值语义；每一轮更新都要保留目标位置或第一个可行值。手算时按这个协议跟踪：拿 banana 手算，重复子串 ana 出现两次，长度 3 可行；若长度 4 不可行，则更长也不可能。然后按协议复查：手算时列出 left、mid、right，以及 mid 处比较结果、可行性检查结果或局部趋势。每一步都要写出被丢弃的一侧为什么不含答案。

\textbf{代码落点。} 检查函数或比较分支单独写清，mid 不溢出，循环退出条件和返回值配套；每个分支都要解释丢弃区间。关键代码行必须能逐句翻译成“旧状态怎样变成新状态”，否则就只是模板。

\textbf{正确性。} 证明从排除一侧开始：答案空间题证明谓词单调，局部比较题证明保留下来的区间仍含目标；不能只靠经验猜测左右边界。

\textbf{边界实验。} 先用全相同字符、无重复、哈希冲突、长度边界、模数选择做反例检查。实验时覆盖全相同字符、无重复、哈希冲突、长度边界、模数选择，逐轮打印 left、mid、right、比较值或检查结果。只要有一次被排除区间仍含答案，二分不变量就错了。

\textbf{复杂度变化。} 暴力逐个试候选是线性或和值域成正比；二分每次排除一半候选，复杂度变成对数级，答案二分还要乘上单次检查成本。

\textbf{迁移追问。} 如果题目改成“也可用后缀数组或后缀自动机，适合系统学习字符串算法”，先判断原状态是否仍然覆盖新答案集合，再决定是微调边界、改 value 语义，还是换算法。

\subsection{LeetCode 480 滑动窗口中位数}

\textbf{题目说明。} LeetCode 480 滑动窗口中位数 的题面入口要先还原成具体任务：窗口移动时既要加入新元素，也要删除滑出的旧元素并查询中位数。

\textbf{样例入口。} 拿 [1, 3, -1]、k=3 手算，中位数是 1；窗口右移时要删除 1、加入 -3，普通堆不能任意删除。

\textbf{暴力基线。} 慢版本每轮扫描全部候选来找最大或最小，再更新候选集合。它的答案选择过程清楚，但动态集合每变化一次就重新找极值，代价太高。放到本题就是：先按“窗口移动时既要加入新元素，也要删除滑出的旧元素并查询中位数”完整试慢版本；样例里已经能看到“规律是双堆维护两半，再用延迟删除清理过期元素，堆顶使用前必须有效”，所以优化前必须先能说清慢版本枚举了哪些候选。

\textbf{暴力过程。} 拿 [1, 3, -1]、k=3 手算，中位数是 1；窗口右移时要删除 1、加入 -3，普通堆不能任意删除。

\textbf{重复工作。} 题面模型：窗口移动时既要加入新元素，也要删除滑出的旧元素并查询中位数。暴力版的慢点要落到本题：慢版本会围绕“窗口移动时既要加入新元素，也要删除滑出的旧元素并查询中位数”反复扫描动态候选集找极值。样例规律是“规律是双堆维护两半，再用延迟删除清理过期元素，堆顶使用前必须有效”，说明每轮真正需要的是堆顶边界，而不是把候选全集重新排序。后面的优化只消掉这类重复工作，不能改变答案集合。

\textbf{优化条件。} 本题的关键观察是：双堆维护左右两半，延迟删除表记录已离窗但尚未到堆顶的元素。这句话必须能翻译成可维护状态；如果题目条件改变后观察不再成立，就不能继续套原来的写法。

\textbf{相邻状态。} 规律是双堆维护两半，再用延迟删除清理过期元素，堆顶使用前必须有效。把这个特例继续拆开看：堆顶必须正好代表下一步最需要处理的候选；若堆里可能有过期对象，读取堆顶前要先清理。堆只解决动态极值，不自动证明被弹出或被丢弃的元素无用。

\textbf{状态复用。} 把每轮全量扫描改成堆顶查询；插入、弹出和过期清理共同维护动态候选集。本题落点是：规律是双堆维护两半，再用延迟删除清理过期元素，堆顶使用前必须有效。手算时只改被新输入影响的那一小部分，而不是回头重算完整候选。

\textbf{指针和字段。} 状态字段要能说清：堆元素字段、堆顶比较规则、候选是否过期、答案读取时机；如果需要懒删除，还要保存删除计数。手算时按这个协议跟踪：拿 [1, 3, -1]、k=3 手算，中位数是 1；窗口右移时要删除 1、加入 -3，普通堆不能任意删除。然后按协议复查：手算时记录堆顶、候选集合和是否有过期元素。每次使用堆顶前都要确认它仍然属于当前问题。

\textbf{代码落点。} 比较器方向先用小样例验证；每次使用堆顶前清理过期项；堆中保存下标时要能回查原数据。关键代码行必须能逐句翻译成“旧状态怎样变成新状态”，否则就只是模板。

\textbf{正确性。} 证明从候选覆盖开始：所有可能成为下一步答案的对象都在堆中，堆顶过期时不会被使用。

\textbf{边界实验。} 先用重复值、偶数窗口、平均值溢出、堆顶过期做反例检查。实验时覆盖重复值、偶数窗口、平均值溢出、堆顶过期，记录堆顶是否过期、比较器是否让正确候选在堆顶、K 或窗口大小为边界值时是否稳定。

\textbf{复杂度变化。} 暴力每轮扫描或排序动态候选，会把线性扫描或排序反复发生；堆把取极值、插入和删除边界控制在对数级。

\textbf{迁移追问。} 如果题目改成“普通 priority\_queue 不能任意删除，所以必须解释懒删除”，先判断原状态是否仍然覆盖新答案集合，再决定是微调边界、改 value 语义，还是换算法。

\subsection{LeetCode 560 和为 K 的子数组}

\textbf{题目说明。} LeetCode 560 和为 K 的子数组给定整数数组 nums 和整数 k，统计和等于 k 的连续子数组个数；数组中可以有负数。

\textbf{样例入口。} 样例 nums=[1, 1, 1]，k=2，输出 2；两个合法连续子数组分别是前两个 1 和后两个 1。

\textbf{暴力基线。} 暴力固定 left，向右累计 sum；每当 sum==k 就计数。

\textbf{暴力过程。} 以 nums=[1, 1, 1]、k=2 为例，从 left=0 可找到 [1, 1]，从 left=1 又找到后两个 1。

\textbf{重复工作。} 题面模型：当前前缀和减去目标值，就是需要匹配的历史前缀。暴力版的慢点要落到本题：浪费在每个右端都回头枚举所有 left。区间和等于 k 可以改写成 prefix[right]-prefix[left]=k。后面的优化只消掉这类重复工作，不能改变答案集合。

\textbf{优化条件。} 本题的关键观察是：哈希表保存前缀和出现次数，先查询再更新。这句话必须能翻译成可维护状态；如果题目条件改变后观察不再成立，就不能继续套原来的写法。

\textbf{相邻状态。} 规律是哈希表保存前缀和频次，且初始要放 0:1。把这个特例继续拆开看：每个合法答案都要还原成当前状态和某个历史状态的配对；哈希表的 key 负责定位历史状态，value 负责表达次数、最早位置或存在性。先查还是先写，必须由是否允许当前前缀和自己配对来决定。

\textbf{状态复用。} 扫描前缀和 prefix。哈希表 count[old\_prefix] 保存历史前缀出现次数；当前位置需要查 prefix-k 出现过几次。手算时只改被新输入影响的那一小部分，而不是回头重算完整候选。

\textbf{指针和字段。} 状态字段要能说清：prefix 是当前前缀和，need=prefix-k 是所需历史前缀，count 保存历史前缀频次，answer 累加匹配次数。手算时按这个协议跟踪：[1, 1, 1] 中扫描到第二个元素时 prefix=2，需要历史 0，count[0]=1，所以找到第一个长度 2 子数组。

\textbf{代码落点。} 关键是先 answer += count[prefix-k]，再 count[prefix]++。初始化 count[0]=1，表示从下标 0 开始的区间。关键代码行必须能逐句翻译成“旧状态怎样变成新状态”，否则就只是模板。

\textbf{正确性。} 任意和为 k 的子数组都对应一对前缀差 k；扫描到右端时，所有合法左端前缀都已在表里。

\textbf{边界实验。} 先用负数、目标为零、从下标零开始的区间、重复前缀做反例检查。实验时把负数、目标为零、从下标零开始的区间、重复前缀加入对拍集合，用平方暴力和优化版比较。失败时先看初始历史状态、查询更新顺序和哈希表 value 类型。

\textbf{复杂度变化。} 暴力 O(\ensuremath{n^{2}})，前缀哈希 O(n) 均摊时间、O(n) 空间。

\textbf{迁移追问。} 如果题目改成“若要求最长区间，哈希表应保存最早位置而不是次数”，先判断原状态是否仍然覆盖新答案集合，再决定是微调边界、改 value 语义，还是换算法。

\subsection{LeetCode 146 LRU 缓存}

\textbf{题目说明。} LeetCode 146 LRU 缓存 的题面入口要先还原成具体任务：哈希表负责按 key 定位，双向链表负责维护新旧顺序。

\textbf{样例入口。} 拿容量 2 依次 put(1)、put(2)、get(1)、put(3) 手算。get(1) 后 1 变成最新，淘汰时应删 2。数组维护顺序每次移动成本高；链表移动节点到头 O(1)，哈希表按 key 定位节点 O(1)。

\textbf{暴力基线。} 暴力用数组保存缓存项，每次 get 线性查找 key，命中后把元素移动到最新位置；put 时线性找 key 或淘汰最旧项。

\textbf{暴力过程。} 容量 2，put(1)、put(2)、get(1)、put(3) 时，get(1) 会让 1 变成最新，put(3) 应淘汰最旧的 2。

\textbf{重复工作。} 题面模型：哈希表负责按 key 定位，双向链表负责维护新旧顺序。暴力版的慢点要落到本题：浪费在按 key 查找和移动新旧顺序都线性。LRU 需要两个能力：O(1) 按 key 定位节点，O(1) 把节点移动到头部和删除尾部。后面的优化只消掉这类重复工作，不能改变答案集合。

\textbf{优化条件。} 本题的关键观察是：访问和更新都把节点移动到头部，容量满时删除尾部最旧节点。这句话必须能翻译成可维护状态；如果题目条件改变后观察不再成立，就不能继续套原来的写法。

\textbf{相邻状态。} 规律是哈希负责找节点，双向链表负责维护新旧顺序。把这个特例继续拆开看：每个指针都要有角色，上一段尾、当前段头、当前节点和后继入口不能混。只要一次改 next 之前没有保存后继，就可能丢掉整段未处理链。

\textbf{状态复用。} 哈希表 key->链表节点迭代器，双向链表按新到旧保存 key/value。访问或更新节点时移动到头部，容量超限删除尾部。手算时只改被新输入影响的那一小部分，而不是回头重算完整候选。

\textbf{指针和字段。} 状态字段要能说清：map 定位缓存项，list.front 是最新项，list.back 是最旧项，capacity 控制淘汰边界。手算时按这个协议跟踪：get(1) 命中后，把 1 对应节点移到表头；put(3) 超容量时，删除表尾 2，并从哈希表同步 erase。

\textbf{代码落点。} C++ 可用 std::list<pair<int, int>> 保存顺序，unordered\_map<int, list::iterator> 定位。更新已有 key 时先改值再 splice 到头部。关键代码行必须能逐句翻译成“旧状态怎样变成新状态”，否则就只是模板。

\textbf{正确性。} 每次访问都把对应项放到最新位置，所以链表尾部始终是最长时间未访问项；哈希表和链表同步维护同一批 key。

\textbf{边界实验。} 先用容量为一、重复 put、get 不存在、更新已有值做反例检查。实验时覆盖容量为一、重复 put、get 不存在、更新已有值，每步画出前驱、当前、后继和下一段头。链表题不要只看输出值，还要检查节点是否丢失或成环。

\textbf{复杂度变化。} 数组暴力 get/put O(n)；哈希表加双链表 get/put 平均 O(1)，空间 O(capacity)。

\textbf{迁移追问。} 如果题目改成“C++ 可用 \code{std::list} 加迭代器，也可手写双向链表理解所有权”，先判断原状态是否仍然覆盖新答案集合，再决定是微调边界、改 value 语义，还是换算法。

\subsection{LeetCode 731 我的日程安排表 II}

\textbf{题目说明。} LeetCode 731 我的日程安排表 II 的题面入口要先还原成具体任务：允许双重预订但不允许三重预订，插入会改变区间重叠层数。

\textbf{样例入口。} 拿 book(10, 20)、book(15, 25) 后，重叠段 [15, 20) 已经是双重预订；再 book(17, 22) 会碰到双重段，必须失败。

\textbf{暴力基线。} 慢版本对新区间和所有旧区间两两比较，手工判断相交、覆盖或冲突。它能验证端点语义，但排序后很多远离当前边界的区间无需再看。放到本题就是：先按“允许双重预订但不允许三重预订，插入会改变区间重叠层数”完整试慢版本；样例里已经能看到“规律是保存所有单次预订和已形成的双重预订，新区间不能与任何双重区间相交”，所以优化前必须先能说清慢版本枚举了哪些候选。

\textbf{暴力过程。} 拿 book(10, 20)、book(15, 25) 后，重叠段 [15, 20) 已经是双重预订；再 book(17, 22) 会碰到双重段，必须失败。

\textbf{重复工作。} 题面模型：允许双重预订但不允许三重预订，插入会改变区间重叠层数。暴力版的慢点要落到本题：慢版本会围绕“允许双重预订但不允许三重预订，插入会改变区间重叠层数”做两两区间比较。样例规律是“规律是保存所有单次预订和已形成的双重预订，新区间不能与任何双重区间相交”，说明排序后只有贴近当前端点、结果尾部或资源堆顶的区间还会影响答案。后面的优化只消掉这类重复工作，不能改变答案集合。

\textbf{优化条件。} 本题的关键观察是：保存已有预订和已经双重预订的区间；新预订若碰到双重区间则失败，否则记录新产生的重叠。这句话必须能翻译成可维护状态；如果题目条件改变后观察不再成立，就不能继续套原来的写法。

\textbf{相邻状态。} 规律是保存所有单次预订和已形成的双重预订，新区间不能与任何双重区间相交。把这个特例继续拆开看：排序以后，真正还能影响当前区间的候选必须贴近当前端点、结果尾部或资源堆顶。某个区间一旦被覆盖、合并或弹出，就要说明它为什么不会和后续区间重新产生更优关系。

\textbf{状态复用。} 把两两比较压成排序后的相邻关系、当前边界或动态有序结构；远离当前边界的区间要被安全归档。本题落点是：规律是保存所有单次预订和已形成的双重预订，新区间不能与任何双重区间相交。手算时只改被新输入影响的那一小部分，而不是回头重算完整候选。

\textbf{指针和字段。} 状态字段要能说清：排序后的区间、当前结果尾、扫描右端或资源堆、端点比较规则；动态版本还要保存前驱后继。手算时按这个协议跟踪：拿 book(10, 20)、book(15, 25) 后，重叠段 [15, 20) 已经是双重预订；再 book(17, 22) 会碰到双重段，必须失败。然后按协议复查：手算时先排序，再逐个区间写当前结果尾部、当前右端或资源堆。每个区间离开视野时都要说明它不会再影响后续。

\textbf{代码落点。} 排序比较器满足严格弱序；端点相等的处理和题意一致；扫描时只和当前结果尾或当前资源集合交互。关键代码行必须能逐句翻译成“旧状态怎样变成新状态”，否则就只是模板。

\textbf{正确性。} 证明从排序后的邻接关系开始：已经合并、选择或丢弃的区间不会被更后面的区间重新影响。

\textbf{边界实验。} 先用端点相接、多个局部重叠、完全覆盖、回滚失败插入做反例检查。实验时覆盖端点相接、多个局部重叠、完全覆盖、回滚失败插入，画出排序后的区间序列和当前结果尾部。动态版本要专门测前驱、后继和端点相接。

\textbf{复杂度变化。} 暴力两两比较区间常见为平方级；排序后扫描通常由排序成本主导，动态有序结构则把每次前驱后继查询控制在对数级。

\textbf{迁移追问。} 如果题目改成“扫描线差分也能做，但要处理失败时恢复状态”，先判断原状态是否仍然覆盖新答案集合，再决定是微调边界、改 value 语义，还是换算法。

\begin{principlebox}{本节复盘方式}
不要只看最终算法名。每道题复盘时先遮住优化版，口述暴力枚举对象；再指出相邻窗口、历史前缀、候选区间、搜索状态或子问题里哪一部分被重复处理；最后用一个边界样例验证状态更新顺序。能讲完这三步，才算真正掌握本章案例。
\end{principlebox}
